\chapter{Closed Localization, Homotopy Purity, and Derived Nilpotent Invariance}
\label{chap:closed-localization-purity-nil-invariance}

\section{Purpose and standing open--closed convention}
\label{sec:closed-localization-purpose}

This chapter develops the closed geometric structure of unstable brave
new motivic homotopy theory.  The preceding chapter associated to every
quasi-compact quasi-separated spectral scheme \(S\) a presentable
Cartesian closed \(\infty\)-category
\[
    \mathbf H(S)
\]
of brave new motivic spaces, together with its pointed smash-monoidal
counterpart
\[
    \mathbf H_\bullet(S).
\]
For every morphism of quasi-compact quasi-separated spectral schemes one
has an adjunction
\[
    f^*
    \colon
    \mathbf H(S)
    \rightleftarrows
    \mathbf H(T)
    :
    f_*,
\]
and for every smooth morphism of finite presentation one has
\[
    p_\sharp
    \colon
    \mathbf H(X)
    \rightleftarrows
    \mathbf H(S)
    :
    p^*.
\]
These adjunctions satisfy smooth base change and the smooth projection
formula, in both unpointed and pointed form.

We now fix a complementary closed--open pair
\[
\begin{tikzcd}[column sep=large]
    Z
        \arrow[r, hook, "i"]
    &
    S
    &
    U
        \arrow[l, hook', "j"']
\end{tikzcd}
\]
in which \(S\) is quasi-compact and quasi-separated,
\[
    i
    \colon
    Z
    \hookrightarrow
    S
\]
is a closed immersion, and
\[
    j
    \colon
    U
    =
    S\setminus Z
    \hookrightarrow
    S
\]
is its quasi-compact open complement.  The quasi-compactness of \(j\)
ensures that \(j\) is smooth and of finite presentation.  Hence the
preceding chapter supplies adjunctions
\[
    i^*
    \colon
    \mathbf H(S)
    \rightleftarrows
    \mathbf H(Z)
    :
    i_*,
\]
and
\[
    j_\sharp
    \colon
    \mathbf H(U)
    \rightleftarrows
    \mathbf H(S)
    :
    j^*
    \rightleftarrows
    \mathbf H(U)
    :
    j_*.
\]
The same notation is used for the corresponding functors on pointed
motivic spaces when no ambiguity can arise.

The central result is the spectral Morel--Voevodsky localization
theorem.  It identifies an arbitrary motivic space over \(S\) as the
gluing of its restriction to \(U\) and its restriction to \(Z\).  From
localization we derive closed recollement, closed base change, the
pointed closed projection formula, pointed smooth--closed exchange, and
derived nilpotent invariance.  We then compare brave new motivic spaces
with ordinary motivic spaces over the classical truncation and transport
ordinary homotopy purity through that comparison to obtain homotopy
purity for smooth closed pairs of spectral schemes.

The chapter remains unstable.  Thom spaces are constructed and
identified geometrically, but their tensor invertibility is reserved
for the stabilization chapter.

All conventions of the preceding chapters remain in force.  In
particular, sheaves are not hypercompleted unless explicitly stated.

\begin{definition}[Motivic spaces supported on a closed subscheme]
\label{def:motivic-spaces-supported-on-closed}
Define
\[
    \mathbf H_Z(S)
    \subseteq
    \mathbf H(S)
\]
to be the full subcategory spanned by the motivic spaces \(F\) such that
\[
    j^*F
    \simeq
    *_U.
\]
Define
\[
    \mathbf H_{\bullet,Z}(S)
    \subseteq
    \mathbf H_\bullet(S)
\]
to be the full subcategory spanned by the pointed motivic spaces \(F\)
such that
\[
    j^*F
    \simeq
    0_U.
\]
\end{definition}

\section{Preliminary open--closed functoriality}
\label{sec:preliminary-open-closed-functoriality}

\subsection{Open immersion functors}

\begin{proposition}[Full faithfulness for open immersions]
\label{prop:open-immersion-full-faithfulness}
Let
\[
    j
    \colon
    U
    \hookrightarrow
    S
\]
be a quasi-compact open immersion.  Then the unit and counit
transformations
\[
    \operatorname{id}_{\mathbf H(U)}
    \longrightarrow
    j^*j_\sharp
\]
and
\[
    j^*j_*
    \longrightarrow
    \operatorname{id}_{\mathbf H(U)}
\]
are equivalences.  Consequently,
\[
    j_\sharp
    \quad\text{and}\quad
    j_*
\]
are fully faithful.  The same assertions hold on pointed motivic
spaces.
\end{proposition}

\begin{proof}
Since \(j\) is a monomorphism, there is a canonical equivalence
\[
    U\times_SU
    \simeq
    U.
\]
Apply smooth base change,
\Ref{thm:unstable-smooth-base-change},
to the Cartesian square
\[
\begin{tikzcd}[column sep=large, row sep=large]
    U
        \arrow[r, "\id_U"]
        \arrow[d, "\id_U"']
    &
    U
        \arrow[d, "j"]
    \\
    U
        \arrow[r, "j"']
    &
    S.
\end{tikzcd}
\]
The resulting equivalence
\[
    j^*j_\sharp
    \simeq
    \operatorname{id}_{\mathbf H(U)}
\]
is the unit of \(j_\sharp\dashv j^*\).  Hence \(j_\sharp\) is fully
faithful.

Let \(F\in\mathbf H(U)\), and let \(X\to U\) be smooth and of finite
presentation.  Since \(j\) is smooth, the composite \(X\to S\) is a
smooth spectral \(S\)-scheme of finite presentation.  Using evaluation
by smooth motives, the adjunction \(j_\sharp\dashv j^*\), smooth direct
image on motives, and the section formula for direct image, we obtain
\[
\begin{aligned}
    \operatorname{Map}_{\mathbf H(U)}
    \bigl(
        \mathrm M_U(X),
        j^*j_*F
    \bigr)
    &\simeq
    \operatorname{Map}_{\mathbf H(S)}
    \bigl(
        j_\sharp\mathrm M_U(X),
        j_*F
    \bigr)
    \\
    &\simeq
    \operatorname{Map}_{\mathbf H(S)}
    \bigl(
        \mathrm M_S(X),
        j_*F
    \bigr)
    \\
    &\simeq
    F(X\times_SU)
    \\
    &\simeq
    F(X).
\end{aligned}
\]
Under these identifications, the map induced by the counit
\[
    j^*j_*F
    \longrightarrow
    F
\]
is the identity on \(F(X)\).  Since smooth motives detect equivalences,
the counit is an equivalence.  Hence \(j_*\) is fully faithful.

The pointed assertions follow from the pointed smooth adjunctions and
the same argument after passage to underlying motivic spaces.
\end{proof}

\subsection{Open--closed orthogonality}

\begin{proposition}[Open--closed orthogonality]
\label{prop:open-closed-orthogonality}
For the complementary pair
\[
    Z
    \xrightarrow{i}
    S
    \xleftarrow{j}
    U,
\]
there are canonical equivalences
\[
    i^*j_\sharp(F)
    \simeq
    \varnothing_Z
\]
for \(F\in\mathbf H(U)\), and
\[
    j^*i_*(G)
    \simeq
    *_U
\]
for \(G\in\mathbf H(Z)\).

On pointed motivic spaces these become
\[
    i^*j_\sharp(F)
    \simeq
    0_Z
\]
and
\[
    j^*i_*(G)
    \simeq
    0_U.
\]
\end{proposition}

\begin{proof}
The fibre product
\[
    U\times_SZ
\]
is empty by the universal property of the open complement established
in
\Ref{sec:closed-immersions-complements}.

Apply smooth base change to the Cartesian square
\[
\begin{tikzcd}[column sep=large, row sep=large]
    \varnothing
        \arrow[r]
        \arrow[d]
    &
    U
        \arrow[d, "j"]
    \\
    Z
        \arrow[r, "i"']
    &
    S.
\end{tikzcd}
\]
Since smooth direct image along
\[
    \varnothing
    \longrightarrow
    Z
\]
carries every object to the initial motivic space, one obtains
\[
    i^*j_\sharp(F)
    \simeq
    \varnothing_Z
\]
for every
\[
    F
    \in
    \mathbf H(U).
\]

For direct image, let
\[
    X
    \longrightarrow
    U
\]
be smooth and of finite presentation.  Then
\[
    X\times_SZ
    \simeq
    \varnothing.
\]
The section formula gives
\[
\begin{aligned}
    (j^*i_*G)(X)
    &\simeq
    (i_*G)(X)
    \\
    &\simeq
    G(X\times_SZ)
    \\
    &\simeq
    G(\varnothing)
    \\
    &\simeq
    *.
\end{aligned}
\]
Hence
\[
    j^*i_*G
    \simeq
    *_U
\]
for every
\[
    G
    \in
    \mathbf H(Z).
\]

We now prove the pointed assertions.  Let
\[
    F
    \in
    \mathbf H_\bullet(U),
\]
and write \(U(F)\) for its underlying unpointed motivic space.  By the
definition of pointed smooth direct image,
\[
    j_{\sharp,\bullet}(F)
    \simeq
    j_\sharp U(F)
    \amalg_{j_\sharp(*_U)}
    *_S.
\]
Since \(i^*\) preserves colimits and terminal objects, the unpointed
orthogonality equivalences give
\[
\begin{aligned}
    i^*j_{\sharp,\bullet}(F)
    &\simeq
    i^*j_\sharp U(F)
    \amalg_{i^*j_\sharp(*_U)}
    *_Z
    \\
    &\simeq
    \varnothing_Z
    \amalg_{\varnothing_Z}
    *_Z
    \\
    &\simeq
    *_Z.
\end{aligned}
\]
With its induced basepoint, this is the zero object
\[
    0_Z
    \in
    \mathbf H_\bullet(Z).
\]
Thus
\[
    i^*j_{\sharp,\bullet}(F)
    \simeq
    0_Z.
\]

Finally, let
\[
    G
    \in
    \mathbf H_\bullet(Z).
\]
The underlying unpointed motivic space of
\[
    j^*i_*G
\]
is
\[
    j^*i_*U(G)
    \simeq
    *_U
\]
by the unpointed assertion.  A pointed motivic space whose underlying
object is terminal is canonically the zero object.  Therefore
\[
    j^*i_*G
    \simeq
    0_U.
\]

We suppress the subscript \(\bullet\) on \(j_\sharp\) in the
statement.
\end{proof}

\subsection{Nisnevich-local conservativity}

\begin{proposition}[Nisnevich-local conservativity]
\label{prop:nisnevich-local-conservativity}
Let
\[
    \{p_\alpha:S_\alpha\to S\}_{\alpha\in A}
\]
be a spectral Nisnevich covering family between quasi-compact
quasi-separated spectral schemes.  Then the pullback functors
\[
    p_\alpha^*
    \colon
    \mathbf H(S)
    \longrightarrow
    \mathbf H(S_\alpha)
\]
are jointly conservative.  The same is true on pointed motivic spaces.
\end{proposition}

\begin{proof}
Let
\[
    \varphi
    \colon
    F
    \longrightarrow
    G
\]
be a morphism in \(\mathbf H(S)\) such that every
\(p_\alpha^*\varphi\) is an equivalence.  Let
\(X\to S\) be smooth and of finite presentation.  The induced family
\[
    X_\alpha
    :=
    X\times_SS_\alpha
    \longrightarrow
    X
\]
is a spectral Nisnevich cover.  Since \(X\) is quasi-compact, a finite
subfamily already covers.  Let
\[
    X_\bullet
    \longrightarrow
    X
\]
be the Čech nerve of that finite subfamily.

Each term \(X_n\) is a finite coproduct of iterated fibre products
\[
    X_{\alpha_0}
    \times_X
    \cdots
    \times_X
    X_{\alpha_n}.
\]
Every such iterated fibre product maps to each of the schemes
\(S_{\alpha_r}\).  Therefore the restriction of \(\varphi\) to every
summand of every \(X_n\) is an equivalence.  Since motivic spaces satisfy
Nisnevich descent and send finite coproducts of smooth schemes to finite
products of spaces, the natural map of totalizations
\[
    \operatorname{Tot}F(X_\bullet)
    \longrightarrow
    \operatorname{Tot}G(X_\bullet)
\]
is an equivalence.  Hence
\[
    F(X)
    \xrightarrow{\simeq}
    G(X).
\]
Since smooth motives detect equivalences by
\Ref{prop:brave-new-motivic-compact-generation},
\(\varphi\) is an equivalence.

The pointed forgetful functor reflects equivalences, so the pointed
assertion follows.
\end{proof}

\subsection{Standard smooth generators}

\begin{definition}[Standard smooth spectral scheme]
\label{def:standard-smooth-spectral-scheme}
A smooth spectral \(S\)-scheme of finite presentation
\[
    X
    \longrightarrow
    S
\]
is \textbf{standard smooth} if \(X\) is affine and there exists an
étale \(S\)-morphism
\[
    X
    \longrightarrow
    \mathbb A^n_S
\]
for some \(n\geq0\).
\end{definition}

\begin{proposition}[Generation by standard smooth motives]
\label{prop:generation-by-standard-smooth-motives}
The category
\[
    \mathbf H(S)
\]
is generated under sifted colimits by the objects
\[
    \mathrm M_S(X)
\]
with \(X\) standard smooth over \(S\).

The pointed category
\[
    \mathbf H_\bullet(S)
\]
is generated under sifted colimits by the objects
\[
    \mathrm M_{S,+}(X)
\]
with \(X\) standard smooth over \(S\).
\end{proposition}

\begin{proof}
By
\cite[Proposition~2.4.4]{KhanBraveNewMotivicI}, in the published version
entitled \emph{The Morel--Voevodsky localization theorem in spectral
algebraic geometry}, the category
\[
    \mathbf H(S)
\]
is generated under sifted colimits by motives
\[
    \mathrm M_S(X)
\]
where \(X\to S\) is affine and smooth and admits an étale
\(S\)-morphism
\[
    X
    \longrightarrow
    \mathbb A_S^n
\]
for some \(n>0\).  Every such \(X\) is standard smooth in the sense of
\Ref{def:standard-smooth-spectral-scheme}.  Hence the larger collection
of all standard smooth motives also generates
\[
    \mathbf H(S)
\]
under sifted colimits.

For the pointed assertion, consider the free-pointing adjunction
\[
    (-)_+
    \colon
    \mathbf H(S)
    \rightleftarrows
    \mathbf H_\bullet(S)
    :
    U_S.
\]
Its associated monad on \(\mathbf H(S)\) is
\[
    T(F)
    =
    F\amalg *_S.
\]
The monadic bar construction expresses every pointed motivic space as
the geometric realization of a simplicial diagram of free pointed
objects.  Geometric realizations are sifted colimits, and the functor
\[
    (-)_+
\]
preserves sifted colimits.

Since every object of \(\mathbf H(S)\) is a sifted colimit of standard
smooth motives, every free pointed motivic space is a sifted colimit of
objects
\[
    \mathrm M_{S,+}(X),
\]
with \(X\) standard smooth.  Applying the bar construction shows that
every object of
\[
    \mathbf H_\bullet(S)
\]
is a sifted colimit of such free pointed motives.
\end{proof}

\section{Exactness of closed direct image}
\label{sec:exactness-closed-direct-image}

\subsection{Local lifting over a closed base}

\begin{proposition}[Local lifting of smooth test schemes]
\label{prop:local-lifting-smooth-test-schemes}
Let
\[
    i
    \colon
    Z
    \hookrightarrow
    S
\]
be a closed immersion, and let
\[
    Y
    \longrightarrow
    Z
\]
be smooth and of finite presentation.  Then every point of \(Y\)
admits a spectral Nisnevich neighbourhood
\[
    Y_\alpha
    \longrightarrow
    Y
\]
for which there exists a smooth spectral \(S\)-scheme of finite
presentation
\[
    \widetilde Y_\alpha
    \longrightarrow
    S
\]
and an equivalence
\[
    \widetilde Y_\alpha\times_SZ
    \simeq
    Y_\alpha.
\]
The same assertion holds with ``smooth'' replaced by ``étale''.
\end{proposition}

\begin{proof}
The local lifting theorem
\Ref{thm:infinitesimal-lifting-smooth-etale}
gives, after shrinking Zariski-locally around each point of \(Y\), a
smooth, respectively étale, lift over \(S\).  A Zariski neighbourhood
is a Nisnevich neighbourhood.  Finite presentation is retained because
the lifted local standard form is a composite of an étale morphism of
finite presentation with brave new affine space of finite presentation.
\end{proof}

\subsection{Weakly contractible colimits}

\begin{definition}[Weakly contractible indexing category]
\label{def:weakly-contractible-indexing-category}
An \(\infty\)-category \(K\) is \textbf{weakly contractible} if its
classifying space
\[
    |K|
\]
is contractible.  A colimit indexed by such a category is called a
\textbf{weakly contractible colimit}.
\end{definition}

\begin{definition}[Topological quasi-cocontinuity]
\label{def:topological-quasi-cocontinuity}
Let \(\mathcal C\) and \(\mathcal D\) be essentially small sites, and
suppose that \(\mathcal D\) admits an initial object
\(\varnothing_{\mathcal D}\).  A functor
\[
    u
    \colon
    \mathcal C
    \longrightarrow
    \mathcal D
\]
is \textbf{topologically quasi-cocontinuous} if, for every object
\(c\in\mathcal C\) and every covering sieve
\[
    R'
    \hookrightarrow
    h_{\mathcal D}(u(c)),
\]
the sieve on \(c\) generated by morphisms \(c'\to c\) for which either
\[
    u(c')
    \simeq
    \varnothing_{\mathcal D}
\]
or the morphism
\[
    h_{\mathcal D}(u(c'))
    \longrightarrow
    h_{\mathcal D}(u(c))
\]
factors through \(R'\), is covering.
\end{definition}

\begin{lemma}[Quasi-cocontinuous restriction]
\label{lem:quasi-cocontinuous-restriction}
Let
\[
    u
    \colon
    \mathcal C
    \longrightarrow
    \mathcal D
\]
be topologically quasi-cocontinuous.  Suppose that the initial object of
\(\mathcal D\) is strict and that the empty sieve covers an object of
\(\mathcal D\) precisely when that object is initial.  Then:
\begin{enumerate}[label=(\roman*)]
    \item restriction
    \[
        u^*
        \colon
        \operatorname{PSh}(\mathcal D)
        \longrightarrow
        \operatorname{PSh}(\mathcal C)
    \]
    sends local equivalences between reduced presheaves to local
    equivalences;

    \item the induced functor on sheaf categories
    \[
        F
        \longmapsto
        L_{\mathcal C}u^*F
    \]
    preserves weakly contractible colimits.
\end{enumerate}
Here a presheaf is reduced if its value on the initial object is
contractible.
\end{lemma}

\begin{proof}
Let
\[
    L_{\mathrm{red}}
    \colon
    \operatorname{PSh}(\mathcal D)
    \longrightarrow
    \operatorname{PSh}_{\mathrm{red}}(\mathcal D)
\]
be reduction.  It forces the value at the initial object to be
contractible and does not change values on noninitial objects.  The
category of reduced presheaves is the free completion of \(\mathcal D\)
under weakly contractible colimits.

Local equivalences in the reduced presheaf category are generated,
under weakly contractible colimits, cobase change, and two-out-of-three,
by covering-sieve morphisms
\[
    R'
    \longrightarrow
    h_{\mathcal D}(d)
\]
with \(d\) noninitial.  Let \(c\in\mathcal C\), and base-change the
restriction of such a morphism along a map
\[
    h_{\mathcal C}(c)
    \longrightarrow
    u^*h_{\mathcal D}(d).
\]
By adjunction and preservation of limits by restriction, the resulting
monomorphism is the sieve generated by maps \(c'\to c\) whose image in
\(\mathcal D\) either becomes initial or factors through the pullback of
\(R'\).  Topological quasi-cocontinuity says precisely that this sieve
is covering.  Thus \(u^*\) sends every generating reduced local
equivalence to a local equivalence, and hence sends every reduced local
equivalence to a local equivalence.  This proves (i).

Every sheaf is reduced under the hypotheses on the initial object.  If
\(K\) is weakly contractible, an objectwise \(K\)-colimit of reduced
presheaves is again reduced because
\[
    \operatorname*{colim}_{k\in K}*
    \simeq
    *.
\]
Applying (i) to the comparison from the objectwise colimit to its
sheafification proves that
\[
    F
    \longmapsto
    L_{\mathcal C}u^*F
\]
preserves \(K\)-colimits.  This proves (ii).  See
\cite{KhanBraveNewMotivicI}.
\end{proof}

\begin{lemma}[Closed base change is quasi-cocontinuous]
\label{lem:closed-base-change-quasi-cocontinuous}
Let
\[
    i
    \colon
    Z
    \hookrightarrow
    S
\]
be a closed immersion.  The base-change functor
\[
    u_i
    \colon
    \operatorname{Sm}^{\mathrm{fp}}_{/S}
    \longrightarrow
    \operatorname{Sm}^{\mathrm{fp}}_{/Z}
\]
is topologically quasi-cocontinuous for the spectral Nisnevich
topologies.
\end{lemma}

\begin{proof}
Let
\[
    X
    \longrightarrow
    S
\]
be smooth and of finite presentation, and let
\[
    R'
    \hookrightarrow
    h_Z(X_Z)
\]
be a Nisnevich covering sieve.  Choose a finite Nisnevich covering
family
\[
    \{V_a\to X_Z\}_{a\in A}
\]
contained in \(R'\).

Consider the closed immersion obtained by base change:
\[
    X_Z
    \hookrightarrow
    X.
\]
For every point
\[
    x\in|X_Z|,
\]
choose an index \(a\in A\) and a point
\[
    v\in|V_a|
\]
above \(x\) which induces an isomorphism of residue fields.  Apply the
local étale lifting theorem
\Ref{prop:local-lifting-smooth-test-schemes}
to the closed immersion
\[
    X_Z
    \hookrightarrow
    X
\]
and to the étale morphism
\[
    V_a
    \longrightarrow
    X_Z.
\]
After replacing \(V_a\) by a Nisnevich neighbourhood
\[
    V_{a,v}
    \longrightarrow
    V_a
\]
of \(v\), there exists an étale morphism of finite presentation
\[
    \widetilde V_{a,v}
    \longrightarrow
    X
\]
and an equivalence
\[
    \widetilde V_{a,v}\times_XX_Z
    \simeq
    V_{a,v}
\]
over \(X_Z\).

The composite
\[
    V_{a,v}
    \longrightarrow
    V_a
    \longrightarrow
    X_Z
\]
belongs to the sieve \(R'\).  Hence the closed base change of
\[
    \widetilde V_{a,v}
    \longrightarrow
    X
\]
factors through \(R'\).  By construction, the lifted family is
Nisnevich-surjective on the points of \(X_Z\).

Adjoin the open immersion
\[
    X_U
    :=
    X\times_SU
    \hookrightarrow
    X,
    \qquad
    U:=S\setminus Z.
\]
The family consisting of
\[
    X_U
    \longrightarrow
    X
\]
and the morphisms
\[
    \widetilde V_{a,v}
    \longrightarrow
    X
\]
is a spectral Nisnevich covering family: the open immersion covers the
complement of \(X_Z\), while the lifted étale morphisms are
Nisnevich-surjective over \(X_Z\).

On \(X_U\), closed base change is empty:
\[
    X_U\times_SZ
    \simeq
    \varnothing.
\]
On every \(\widetilde V_{a,v}\), closed base change factors through
\(R'\).  Therefore the sieve on \(X\) generated by morphisms
\[
    X'\longrightarrow X
\]
for which either
\[
    X'_Z
    \simeq
    \varnothing
\]
or
\[
    h_Z(X'_Z)
    \longrightarrow
    h_Z(X_Z)
\]
factors through \(R'\), is covering.  This is precisely the
topological quasi-cocontinuity of \(u_i\).
\end{proof}

\begin{lemma}[Motivic quasi-cocontinuous restriction]
\label{lem:motivic-quasi-cocontinuous-restriction}
Let
\[
    u
    \colon
    \mathcal C
    \longrightarrow
    \mathcal D
\]
be a topologically quasi-cocontinuous functor between essentially small
sites with strict initial objects.  Suppose that:
\begin{enumerate}[label=(\roman*)]
    \item the empty sieve covers precisely the initial object;

    \item
    \[
        u(\varnothing_{\mathcal C})
        \simeq
        \varnothing_{\mathcal D};
    \]

    \item both sites are equipped with affine-line objects
    \(\mathbb A^1_{\mathcal C}\) and
    \(\mathbb A^1_{\mathcal D}\), together with their zero and unit
    sections and scalar multiplication;

    \item the functor \(u\) preserves terminal objects and there are
    natural equivalences
    \[
        u(c\times\mathbb A^1_{\mathcal C})
        \simeq
        u(c)\times\mathbb A^1_{\mathcal D}
    \]
    compatible with the projections, zero and unit sections, and scalar
    multiplication.
\end{enumerate}
Then restriction
\[
    u^*
    \colon
    \operatorname{PSh}(\mathcal D)
    \longrightarrow
    \operatorname{PSh}(\mathcal C)
\]
sends motivic equivalences between reduced presheaves to motivic
equivalences.
\end{lemma}

\begin{proof}
By
\Ref{lem:quasi-cocontinuous-restriction},
restriction sends Nisnevich-local equivalences between reduced
presheaves to Nisnevich-local equivalences.

Let \(d\in\mathcal D\), and consider the affine-line projection
\[
    p_d
    \colon
    h_{\mathcal D}
    \bigl(
        d\times\mathbb A^1_{\mathcal D}
    \bigr)
    \longrightarrow
    h_{\mathcal D}(d).
\]
The zero section determines a morphism
\[
    z_d
    \colon
    h_{\mathcal D}(d)
    \longrightarrow
    h_{\mathcal D}
    \bigl(
        d\times\mathbb A^1_{\mathcal D}
    \bigr)
\]
satisfying
\[
    p_dz_d
    =
    \operatorname{id}.
\]
Scalar multiplication on the affine-line factor defines an
\(\mathbb A^1_{\mathcal D}\)-homotopy from the identity of
\[
    h_{\mathcal D}
    \bigl(
        d\times\mathbb A^1_{\mathcal D}
    \bigr)
\]
to
\[
    z_dp_d.
\]
Condition (iv), applied to the terminal object of \(\mathcal C\),
gives an equivalence
\[
    u(\mathbb A^1_{\mathcal C})
    \simeq
    \mathbb A^1_{\mathcal D}.
\]
Functoriality of \(u\) therefore gives a natural morphism
\[
    h_{\mathcal C}(\mathbb A^1_{\mathcal C})
    \longrightarrow
    u^*h_{\mathcal D}(\mathbb A^1_{\mathcal D}).
\]
Restricting the preceding homotopy and precomposing its parameter with
this morphism gives an \(\mathbb A^1_{\mathcal C}\)-homotopy from the
identity of
\[
    u^*h_{\mathcal D}
    \bigl(
        d\times\mathbb A^1_{\mathcal D}
    \bigr)
\]
to
\[
    u^*z_d\,u^*p_d.
\]
Thus
\[
    u^*p_d
\]
is an elementary \(\mathbb A^1_{\mathcal C}\)-homotopy equivalence,
with homotopy inverse \(u^*z_d\).  In particular, it is an
\(\mathbb A^1_{\mathcal C}\)-local equivalence.

The motivic equivalences between reduced presheaves are generated from
the reduced Nisnevich-local equivalences and the affine-line
projections by weakly contractible colimits, cobase change, retracts,
and two-out-of-three.  Restriction preserves these operations.  The
preceding two paragraphs therefore show that restriction carries every
motivic equivalence between reduced presheaves to a motivic
equivalence.

See
\cite{KhanBraveNewMotivicI}.
\end{proof}

\begin{lemma}[Closed restriction preserves motivic-local objects]
\label{lem:closed-restriction-preserves-motivic-local-objects}
Let
\[
    i
    \colon
    Z
    \hookrightarrow
    S
\]
be a closed immersion, and let
\[
    u_i
    \colon
    \operatorname{Sm}^{\mathrm{fp}}_{/S}
    \longrightarrow
    \operatorname{Sm}^{\mathrm{fp}}_{/Z},
    \qquad
    X
    \longmapsto
    X_Z,
\]
be closed base change.  If
\[
    F
    \in
    \operatorname{PSh}(Z)
\]
is motivic-local, then
\[
    u_i^*F
    \in
    \operatorname{PSh}(S)
\]
is motivic-local.  On motivic-local presheaves, restriction along
\(u_i\) is the direct image functor
\[
    i_*
    \colon
    \mathbf H(Z)
    \longrightarrow
    \mathbf H(S).
\]
\end{lemma}

\begin{proof}
Since
\[
    \varnothing_S\times_SZ
    \simeq
    \varnothing_Z,
\]
one has
\[
    (u_i^*F)(\varnothing_S)
    =
    F(\varnothing_Z)
    \simeq
    *.
\]

Let
\[
\begin{tikzcd}[column sep=large, row sep=large]
    W \arrow[r] \arrow[d]
    &
    V \arrow[d]
    \\
    U \arrow[r]
    &
    X
\end{tikzcd}
\]
be a spectral Nisnevich distinguished square over \(S\).  Closed base
change preserves open immersions, étale morphisms, fibre products, and
the equivalence over the closed complement.  Hence
\[
\begin{tikzcd}[column sep=large, row sep=large]
    W_Z \arrow[r] \arrow[d]
    &
    V_Z \arrow[d]
    \\
    U_Z \arrow[r]
    &
    X_Z
\end{tikzcd}
\]
is a spectral Nisnevich distinguished square over \(Z\).  Since \(F\)
is Nisnevich-local, the canonical morphism
\[
    F(X_Z)
    \longrightarrow
    F(U_Z)
    \times_{F(W_Z)}
    F(V_Z)
\]
is an equivalence.  This is the Nisnevich-excision condition for
\(u_i^*F\).

Finally, closed base change is compatible with the brave new affine
line:
\[
    (X\times_S\mathbb A^1_S)_Z
    \simeq
    X_Z\times_Z\mathbb A^1_Z.
\]
Therefore
\[
\begin{aligned}
    (u_i^*F)(X)
    &=
    F(X_Z)
    \\
    &\simeq
    F(X_Z\times_Z\mathbb A^1_Z)
    \\
    &\simeq
    (u_i^*F)(X\times_S\mathbb A^1_S).
\end{aligned}
\]
Thus \(u_i^*F\) is motivic-local.  The final assertion is the section
formula defining motivic direct image.
\end{proof}

\begin{theorem}[Closed direct image preserves weakly contractible colimits]
\label{thm:closed-direct-image-weakly-contractible-colimits}
Let
\[
    i
    \colon
    Z
    \hookrightarrow
    S
\]
be a closed immersion between quasi-compact quasi-separated spectral
schemes, with quasi-compact open complement.  Then
\[
    i_*
    \colon
    \mathbf H(Z)
    \longrightarrow
    \mathbf H(S)
\]
preserves weakly contractible colimits.
\end{theorem}

\begin{proof}
At the presheaf level, direct image is restriction along the closed
base-change functor
\[
    u_i
    \colon
    \operatorname{Sm}^{\mathrm{fp}}_{/S}
    \longrightarrow
    \operatorname{Sm}^{\mathrm{fp}}_{/Z}.
\]
By
\Ref{lem:closed-base-change-quasi-cocontinuous},
this functor is topologically quasi-cocontinuous.  It preserves the
initial and terminal objects.  Closed base change is compatible with
the brave new affine line and all its structure maps:
\[
    (X\times_S\mathbb A^1_S)_Z
    \simeq
    X_Z\times_Z\mathbb A^1_Z.
\]
Thus the hypotheses of
\Ref{lem:motivic-quasi-cocontinuous-restriction}
are satisfied.

Let \(K\) be weakly contractible, and let
\[
    F_\bullet
    \colon
    K
    \longrightarrow
    \mathbf H(Z)
\]
be a diagram.  Regard the \(F_k\) as their underlying motivic-local
presheaves, and put
\[
    G
    :=
    \operatorname*{colim}_{k\in K}^{\operatorname{PSh}(Z)}F_k.
\]
Since each \(F_k\) is reduced and \(K\) is weakly contractible,
\[
    G(\varnothing)
    \simeq
    \operatorname*{colim}_{k\in K}*
    \simeq
    *.
\]
Hence \(G\) is reduced.  The colimit of the diagram in
\(\mathbf H(Z)\) is
\[
    \operatorname*{colim}_{k\in K}^{\mathbf H(Z)}F_k
    \simeq
    L_{\mathrm{mot},Z}G.
\]

By
\Ref{lem:closed-restriction-preserves-motivic-local-objects},
restriction along \(u_i\) carries motivic-local presheaves over \(Z\)
to motivic-local presheaves over \(S\), and on such presheaves it is the
direct image \(i_*\).  Since restriction computes colimits objectwise,
\[
\begin{aligned}
    \operatorname*{colim}_{k\in K}^{\mathbf H(S)}i_*F_k
    &\simeq
    L_{\mathrm{mot},S}
    \operatorname*{colim}_{k\in K}^{\operatorname{PSh}(S)}u_i^*F_k
    \\
    &\simeq
    L_{\mathrm{mot},S}u_i^*G.
\end{aligned}
\]
The localization unit
\[
    G
    \longrightarrow
    L_{\mathrm{mot},Z}G
\]
is a motivic equivalence between reduced presheaves.  By
\Ref{lem:motivic-quasi-cocontinuous-restriction}, its restriction
\[
    u_i^*G
    \longrightarrow
    u_i^*L_{\mathrm{mot},Z}G
\]
is a motivic equivalence over \(S\).  Its target is motivic-local, so
localizing the source gives an equivalence
\[
    L_{\mathrm{mot},S}u_i^*G
    \xrightarrow{\simeq}
    u_i^*L_{\mathrm{mot},Z}G.
\]
The right side is
\[
    i_*
    \operatorname*{colim}_{k\in K}^{\mathbf H(Z)}F_k.
\]
Combining the preceding equivalences gives
\[
    \operatorname*{colim}_{k\in K}i_*F_k
    \xrightarrow{\simeq}
    i_*
    \operatorname*{colim}_{k\in K}F_k.
\]
Thus \(i_*\) preserves weakly contractible colimits.
\end{proof}

\begin{corollary}[Closed restriction and motivic localization]
\label{cor:closed-restriction-motivic-localization}
Let
\[
    G
    \in
    \operatorname{PSh}(Z)
\]
be reduced.  Then there is a canonical equivalence
\[
    L_{\mathrm{mot},S}
    i_{*,\operatorname{PSh}}G
    \simeq
    i_*L_{\mathrm{mot},Z}G.
\]
\end{corollary}

\begin{proof}
The localization unit
\[
    G
    \longrightarrow
    L_{\mathrm{mot},Z}G
\]
is a motivic equivalence.  Its source is reduced by assumption, and its
target is reduced because every motivic-local presheaf takes the
initial object to a contractible space.

By
\Ref{lem:motivic-quasi-cocontinuous-restriction},
restriction along the closed base-change functor carries this
localization unit to a motivic equivalence
\[
    i_{*,\operatorname{PSh}}G
    \longrightarrow
    i_{*,\operatorname{PSh}}
    L_{\mathrm{mot},Z}G
\]
over \(S\).

By
\Ref{lem:closed-restriction-preserves-motivic-local-objects},
the target is motivic-local and agrees with
\[
    i_*L_{\mathrm{mot},Z}G.
\]
Applying motivic localization to the source therefore gives a canonical
equivalence
\[
    L_{\mathrm{mot},S}
    i_{*,\operatorname{PSh}}G
    \xrightarrow{\simeq}
    i_*L_{\mathrm{mot},Z}G.
\]
\end{proof}

\begin{lemma}[Pointed colimits from weakly contractible colimits]
\label{lem:pointed-colimits-from-weakly-contractible}
Let
\[
    R
    \colon
    \mathcal C
    \longrightarrow
    \mathcal D
\]
be an accessible functor between presentable \(\infty\)-categories
which preserves terminal objects and weakly contractible colimits.
Then the induced functor
\[
    R_\bullet
    \colon
    \mathcal C_{*/}
    \longrightarrow
    \mathcal D_{*/}
\]
preserves all small colimits.
\end{lemma}

\begin{proof}
Let \(D:K\to\mathcal C_{*/}\) be a small diagram.  Adjoin an initial
cone point to \(K\), obtaining \(K^\triangleleft\).  Extend the
underlying diagram to \(K^\triangleleft\) by sending the cone point to
the terminal object of \(\mathcal C\) and the cone arrows to the chosen
basepoint maps.  The colimit of this extended diagram in \(\mathcal C\)
is the underlying object of the colimit of \(D\) in
\(\mathcal C_{*/}\).  Since \(K^\triangleleft\) has an initial object,
it is weakly contractible.  The hypotheses on \(R\) therefore identify
the image of this colimit with the corresponding colimit in
\(\mathcal D_{*/}\).  Thus \(R_\bullet\) preserves the colimit of every
small diagram.
\end{proof}

\begin{corollary}[Colimit preservation of pointed closed direct image]
\label{cor:pointed-closed-direct-image-colimit-preserving}
The pointed closed direct image
\[
    i_*
    \colon
    \mathbf H_\bullet(Z)
    \longrightarrow
    \mathbf H_\bullet(S)
\]
preserves all small colimits.
\end{corollary}

\begin{proof}
The unpointed direct image preserves terminal objects because it is a
right adjoint, and it preserves weakly contractible colimits by
\Ref{thm:closed-direct-image-weakly-contractible-colimits}.  Apply
\Ref{lem:pointed-colimits-from-weakly-contractible}.
\end{proof}

\subsection{Exceptional inverse image for closed immersions}

\begin{proposition}[Exceptional inverse image for a closed immersion]
\label{prop:closed-exceptional-inverse-image}
The pointed closed direct image
\[
    i_*
    \colon
    \mathbf H_\bullet(Z)
    \longrightarrow
    \mathbf H_\bullet(S)
\]
admits a right adjoint
\[
    i^!
    \colon
    \mathbf H_\bullet(S)
    \longrightarrow
    \mathbf H_\bullet(Z).
\]
Thus there is an adjunction
\[
    i_*
    \dashv
    i^!.
\]
\end{proposition}

\begin{proof}
By
\Ref{cor:pointed-closed-direct-image-colimit-preserving},
the functor \(i_*\) preserves all small colimits.  It is accessible
because it is induced by an accessible functor between the underlying
presentable motivic categories.

Both
\[
    \mathbf H_\bullet(Z)
    \quad\text{and}\quad
    \mathbf H_\bullet(S)
\]
are presentable.  The adjoint functor theorem therefore supplies a
right adjoint
\[
    i^!.
\]
\end{proof}

\section{The spectral Morel--Voevodsky localization theorem}
\label{sec:spectral-morel-voevodsky-localization}

\subsection{The localization transformation}

\begin{construction}[The localization square]
\label{con:localization-square}
Let
\[
    F
    \in
    \mathbf H(S).
\]
The counit of
\[
    j_\sharp
    \dashv
    j^*
\]
gives
\[
    j_\sharp j^*F
    \longrightarrow
    F,
\]
and the unit of
\[
    i^*
    \dashv
    i_*
\]
gives
\[
    F
    \longrightarrow
    i_*i^*F.
\]
By
\Ref{prop:open-closed-orthogonality},
their composite factors canonically through the unique morphism
\[
    j_\sharp j^*F
    \longrightarrow
    j_\sharp(*_U).
\]
The terminal map \(j^*F\to *_U\) gives a commutative square
\[
\begin{tikzcd}[column sep=large, row sep=large]
    j_\sharp j^*F
        \arrow[r]
        \arrow[d]
    &
    F
        \arrow[d]
    \\
    j_\sharp(*_U)
        \arrow[r]
    &
    i_*i^*F.
\end{tikzcd}
\]
This is the \textbf{localization square} of \(F\).
\end{construction}

\begin{theorem}[Spectral Morel--Voevodsky localization]
\label{thm:spectral-morel-voevodsky-localization}
For every
\[
    F
    \in
    \mathbf H(S),
\]
the localization square
\[
\begin{tikzcd}[column sep=large, row sep=large]
    j_\sharp j^*F
        \arrow[r]
        \arrow[d]
    &
    F
        \arrow[d]
    \\
    j_\sharp(*_U)
        \arrow[r]
    &
    i_*i^*F
\end{tikzcd}
\]
is cocartesian.
\end{theorem}

The proof occupies the remainder of this section.

\subsection{Presheaves supported near the closed subscheme}

\begin{construction}[The presheaf obtained by collapsing the open part]
\label{con:presheaf-collapsing-open-part}
Let
\[
    X
    \longrightarrow
    S
\]
be smooth and of finite presentation.  Put
\[
    X_U
    :=
    X\times_SU,
    \qquad
    X_Z
    :=
    X\times_SZ.
\]
Define
\[
    h_S^Z(X)
    :=
    h_S(X)
    \amalg_{h_S(X_U)}
    h_S(U).
\]
At the presheaf level, restriction along closed base change gives a
functor
\[
    i_{*,\operatorname{PSh}}
    \colon
    \operatorname{PSh}(Z)
    \longrightarrow
    \operatorname{PSh}(S).
\]
There is a canonical morphism
\[
    \rho_X
    \colon
    h_S^Z(X)
    \longrightarrow
    i_{*,\operatorname{PSh}}h_Z(X_Z).
\]
On the summand \(h_S(X)\), it is induced by restriction to \(Z\).  On
the summand \(h_S(U)\), it is induced by the empty closed base change.
\end{construction}

\begin{proposition}[Localization square on a smooth motive]
\label{prop:localization-square-smooth-motive}
For
\[
    X
    \in
    \operatorname{Sm}^{\mathrm{fp}}_{/S},
\]
the upper-left pushout in the localization square of
\(\mathrm M_S(X)\) is canonically equivalent to
\[
    L_{\mathrm{mot},S}h_S^Z(X).
\]
Under this identification, the induced map to
\[
    i_*i^*\mathrm M_S(X)
    \simeq
    i_*\mathrm M_Z(X_Z)
\]
is obtained by applying motivic localization to
\[
    \rho_X
    \colon
    h_S^Z(X)
    \longrightarrow
    i_{*,\operatorname{PSh}}h_Z(X_Z)
\]
and then using the canonical equivalence
\[
    L_{\mathrm{mot},S}
    i_{*,\operatorname{PSh}}h_Z(X_Z)
    \simeq
    i_*\mathrm M_Z(X_Z).
\]
\end{proposition}

\begin{proof}
By
\Ref{thm:smooth-motivic-direct-image},
\[
    j_\sharp\mathrm M_U(X_U)
    \simeq
    \mathrm M_S(X_U),
\]
where \(X_U\) is regarded over \(S\) by composition.  Moreover,
\[
    j_\sharp(*_U)
    \simeq
    \mathrm M_S(U),
\]
since
\[
    *_U
    =
    \mathrm M_U(U).
\]
Since motivic localization preserves pushouts, the pushout
\[
    \mathrm M_S(X)
    \amalg_{\mathrm M_S(X_U)}
    \mathrm M_S(U)
\]
is canonically equivalent to
\[
    L_{\mathrm{mot},S}
    \left(
        h_S(X)
        \amalg_{h_S(X_U)}
        h_S(U)
    \right)
    =
    L_{\mathrm{mot},S}h_S^Z(X).
\]

Base change gives
\[
    i^*\mathrm M_S(X)
    \simeq
    \mathrm M_Z(X_Z).
\]
The representable presheaf
\[
    h_Z(X_Z)
\]
is reduced, because
\[
    h_Z(X_Z)(\varnothing_Z)
    \simeq
    *.
\]
Therefore
\Ref{cor:closed-restriction-motivic-localization}
gives a canonical equivalence
\[
\begin{aligned}
    L_{\mathrm{mot},S}
    i_{*,\operatorname{PSh}}h_Z(X_Z)
    &\simeq
    i_*L_{\mathrm{mot},Z}h_Z(X_Z)
    \\
    &=
    i_*\mathrm M_Z(X_Z).
\end{aligned}
\]

Under these identifications, the morphism from the upper-left pushout
to
\[
    i_*i^*\mathrm M_S(X)
    \simeq
    i_*\mathrm M_Z(X_Z)
\]
is precisely the composite
\[
\begin{aligned}
    L_{\mathrm{mot},S}h_S^Z(X)
    &\xrightarrow{
        L_{\mathrm{mot},S}(\rho_X)
    }
    L_{\mathrm{mot},S}
    i_{*,\operatorname{PSh}}h_Z(X_Z)
    \\
    &\xrightarrow{\simeq}
    i_*\mathrm M_Z(X_Z).
\end{aligned}
\]
\end{proof}

\subsection{Trivialized maps}

\begin{construction}[Presheaf of \(Z\)-trivialized maps]
\label{con:presheaf-z-trivialized-maps}
Let
\[
    X
    \longrightarrow
    S
\]
be smooth and of finite presentation, and let
\[
    t
    \colon
    Z
    \longrightarrow
    X
\]
be an \(S\)-morphism.  Equivalently, \(t\) is a section of
\[
    X_Z
    \longrightarrow
    Z.
\]
The section \(t\) determines a morphism
\[
    *_S
    \longrightarrow
    i_{*,\operatorname{PSh}}h_Z(X_Z).
\]
Define
\[
    h_S(X,t)
    :=
    h_S^Z(X)
    \times_{i_{*,\operatorname{PSh}}h_Z(X_Z)}
    *_S.
\]
\end{construction}

A section of \(h_S(X,t)\) over a smooth spectral \(S\)-scheme
\(Y\) consists of an \(S\)-morphism
\[
    f
    \colon
    Y
    \longrightarrow
    X
\]
together with an identification of its restriction to
\[
    Y_Z
    :=
    Y\times_SZ
\]
with the composite
\[
    Y_Z
    \longrightarrow
    Z
    \xrightarrow{t}
    X.
\]
If \(Y_Z\) is empty, the value is canonically contractible.

\begin{proposition}[Smooth base change of trivialized-map presheaves]
\label{prop:trivialized-map-base-change}
Let
\[
    f
    \colon
    T
    \longrightarrow
    S
\]
be a smooth morphism of finite presentation.  Form the Cartesian square
\[
\begin{tikzcd}[column sep=large, row sep=large]
    Z_T
        \arrow[r, "f_Z"]
        \arrow[d, "i_T"']
    &
    Z
        \arrow[d, "i"]
    \\
    T
        \arrow[r, "f"']
    &
    S.
\end{tikzcd}
\]
Put
\[
    X_T
    :=
    X\times_ST,
\]
and let
\[
    t_T
    \colon
    Z_T
    \longrightarrow
    X_T
\]
be the base change of \(t\).

There is a canonical presheaf-level exchange equivalence
\[
    f^*_{\operatorname{PSh}}
    i_{*,\operatorname{PSh}}
    \simeq
    i_{T,*,\operatorname{PSh}}
    f_Z^*.
\]
Consequently, there is a canonical equivalence
\[
    f^*_{\operatorname{PSh}}h_S(X,t)
    \simeq
    h_T(X_T,t_T).
\]
After motivic localization, one obtains
\[
    f^*L_{\mathrm{mot},S}h_S(X,t)
    \simeq
    L_{\mathrm{mot},T}h_T(X_T,t_T).
\]
\end{proposition}

\begin{proof}
Because \(f\) is smooth and of finite presentation, presheaf inverse
image is restriction along the composition functor
\[
    f^{\mathrm{geom}}_\sharp
    \colon
    \mathcal C_T
    \longrightarrow
    \mathcal C_S.
\]
It therefore preserves all limits and colimits.

Let
\[
    H
    \in
    \operatorname{PSh}(Z),
\]
and let
\[
    Y
    \in
    \mathcal C_T.
\]
There are natural equivalences
\[
\begin{aligned}
    \bigl(
        f^*_{\operatorname{PSh}}
        i_{*,\operatorname{PSh}}H
    \bigr)(Y)
    &\simeq
    \bigl(
        i_{*,\operatorname{PSh}}H
    \bigr)(Y)
    \\
    &\simeq
    H(Y\times_SZ),
\end{aligned}
\]
where \(Y\) in the middle expression is regarded as an \(S\)-scheme by
composition with \(f\).  On the other hand,
\[
\begin{aligned}
    \bigl(
        i_{T,*,\operatorname{PSh}}
        f_Z^*H
    \bigr)(Y)
    &\simeq
    \bigl(
        f_Z^*H
    \bigr)(Y\times_TZ_T)
    \\
    &\simeq
    H(Y\times_TZ_T)
    \\
    &\simeq
    H(Y\times_SZ).
\end{aligned}
\]
These equivalences are natural in \(H\) and \(Y\), and hence determine
the canonical exchange equivalence
\[
    f^*_{\operatorname{PSh}}
    i_{*,\operatorname{PSh}}
    \simeq
    i_{T,*,\operatorname{PSh}}
    f_Z^*.
\]

Smooth inverse image sends representables to their geometric base
changes:
\[
\begin{aligned}
    f^*_{\operatorname{PSh}}h_S(X)
    &\simeq
    h_T(X_T),
    \\
    f^*_{\operatorname{PSh}}h_S(X_U)
    &\simeq
    h_T((X_T)_{U_T}),
    \\
    f^*_{\operatorname{PSh}}h_S(U)
    &\simeq
    h_T(U_T).
\end{aligned}
\]
Since it preserves pushouts, it follows that
\[
    f^*_{\operatorname{PSh}}h_S^Z(X)
    \simeq
    h_T^{Z_T}(X_T).
\]

The exchange equivalence gives
\[
\begin{aligned}
    f^*_{\operatorname{PSh}}
    \bigl(
        i_{*,\operatorname{PSh}}h_Z(X_Z)
    \bigr)
    &\simeq
    i_{T,*,\operatorname{PSh}}
    f_Z^*h_Z(X_Z)
    \\
    &\simeq
    i_{T,*,\operatorname{PSh}}
    h_{Z_T}(X_{Z_T}).
\end{aligned}
\]
The point determined by \(t\) is carried to the point determined by
\(t_T\).  Applying \(f^*_{\operatorname{PSh}}\) to the pullback
defining \(h_S(X,t)\) therefore gives
\[
\begin{aligned}
    f^*_{\operatorname{PSh}}h_S(X,t)
    &\simeq
    h_T^{Z_T}(X_T)
    \times_{
        i_{T,*,\operatorname{PSh}}
        h_{Z_T}(X_{Z_T})
    }
    *_T
    \\
    &\simeq
    h_T(X_T,t_T).
\end{aligned}
\]

Compatibility of motivic inverse image with motivic localization gives
the final equivalence.
\end{proof}

\subsection{Étale invariance of trivialized maps}

\begin{proposition}[Étale invariance of trivialized-map presheaves]
\label{prop:etale-invariance-trivialized-maps}
Let
\[
    X
    \longrightarrow
    S
\]
be smooth and of finite presentation.  Consider an étale morphism over
\(S\)
\[
    q
    \colon
    X'
    \longrightarrow
    X
\]
and \(S\)-morphisms
\[
    t'
    \colon
    Z
    \longrightarrow
    X',
    \qquad
    t
    \colon
    Z
    \longrightarrow
    X
\]
such that
\[
    q\circ t'
    =
    t.
\]
Write
\[
    q_Z
    \colon
    X'_Z
    \longrightarrow
    X_Z
\]
for the base change of \(q\).  Suppose that the induced square
\[
\begin{tikzcd}[column sep=large, row sep=large]
    Z
        \arrow[r, "t'"]
        \arrow[d, "\id_Z"']
    &
    X'_Z
        \arrow[d, "q_Z"]
    \\
    Z
        \arrow[r, "t"']
    &
    X_Z
\end{tikzcd}
\]
is Cartesian.  Then
\[
    h_S(X',t')
    \longrightarrow
    h_S(X,t)
\]
is a Nisnevich-local equivalence and hence a motivic equivalence.
\end{proposition}

\begin{proof}
Since \(q\) is étale, it is smooth and of finite presentation.
Consequently, the composite
\[
    X'
    \longrightarrow
    X
    \longrightarrow
    S
\]
is smooth and of finite presentation, so both trivialized-map
presheaves in the statement are defined.

Write
\[
    \varphi
    \colon
    h_S(X',t')
    \longrightarrow
    h_S(X,t).
\]
We prove that the morphism obtained after Nisnevich sheafification is
an equivalence.

We first show that
\[
    L_{\mathrm{Nis},S}(\varphi)
\]
is \(0\)-truncated.  Since Nisnevich sheafification is left exact, it
is enough to prove that \(\varphi\) is already \(0\)-truncated as a
morphism of presheaves.

Let
\[
    Y
    \longrightarrow
    S
\]
be smooth and of finite presentation.  If
\[
    Y_Z
    :=
    Y\times_SZ
    \simeq
    \varnothing,
\]
then both
\[
    h_S(X',t')(Y)
    \qquad\text{and}\qquad
    h_S(X,t)(Y)
\]
are canonically contractible.  Hence
\[
    \varphi(Y)
\]
is an equivalence and in particular is \(0\)-truncated.

Suppose now that
\[
    Y_Z
    \not\simeq
    \varnothing.
\]
In this case evaluation of \(\varphi\) is the morphism induced on
fibres in the diagram
\[
\begin{tikzcd}[column sep=large]
    h_S(X',t')(Y)
        \arrow[r]
        \arrow[d]
    &
    \operatorname{Map}_S(Y,X')
        \arrow[r]
        \arrow[d]
    &
    \operatorname{Map}_Z(Y_Z,X'_Z)
        \arrow[d]
    \\
    h_S(X,t)(Y)
        \arrow[r]
    &
    \operatorname{Map}_S(Y,X)
        \arrow[r]
    &
    \operatorname{Map}_Z(Y_Z,X_Z).
\end{tikzcd}
\]
The two right-hand vertical morphisms are \(0\)-truncated because
\(q\) and \(q_Z\) are étale.  It follows that the left-hand vertical
morphism is \(0\)-truncated.  Thus \(\varphi\) is objectwise
\(0\)-truncated, and consequently
\[
    L_{\mathrm{Nis},S}(\varphi)
\]
is \(0\)-truncated.

We next show that
\[
    L_{\mathrm{Nis},S}(\varphi)
\]
is an effective epimorphism.  Let
\[
    (f,\eta)
    \in
    h_S(X,t)(Y)
\]
be a section.  The case \(Y_Z\simeq\varnothing\) is immediate, since
both evaluated presheaves are contractible.  We may therefore suppose
that \(Y_Z\) is nonempty.

Pulling \(q\) back along \(f\) gives an étale morphism
\[
    q_Y
    \colon
    Y'
    :=
    Y\times_XX'
    \longrightarrow
    Y.
\]
The trivialization \(\eta\) identifies the restriction of \(f\) to
\(Y_Z\) with the composite
\[
    Y_Z
    \longrightarrow
    Z
    \xrightarrow{t}
    X_Z.
\]
Therefore
\[
\begin{aligned}
    Y'\times_YY_Z
    &\simeq
    Y_Z\times_{X_Z}X'_Z
    \\
    &\simeq
    Y_Z\times_Z
    \bigl(
        Z\times_{X_Z}X'_Z
    \bigr)
    \\
    &\simeq
    Y_Z,
\end{aligned}
\]
where the last equivalence uses the Cartesian square in the
statement.

Consequently,
\[
    \{Y_U\to Y,\;Y'\to Y\},
    \qquad
    Y_U:=Y\times_SU,
\]
is a spectral Nisnevich covering family.  Over \(Y'\), the morphism
\(f\) has the tautological lift to \(X'\).  Over \(Y_U\), the closed
base change is empty, so both trivialized-map spaces are contractible.
Thus every section of the target lifts Nisnevich-locally.  Hence
\[
    L_{\mathrm{Nis},S}(\varphi)
\]
is an effective epimorphism.

It remains to show that the diagonal of
\[
    L_{\mathrm{Nis},S}(\varphi)
\]
is an effective epimorphism.  Let two sections
\[
    (f',\eta_{f'}),
    \qquad
    (g',\eta_{g'})
    \in
    h_S(X',t')(Y)
\]
be equipped with an identification of their images in
\[
    h_S(X,t)(Y).
\]
Again, if \(Y_Z\simeq\varnothing\), the claim is immediate because the
relevant spaces are contractible.  We may therefore suppose that
\(Y_Z\) is nonempty.

Since the diagonal of the étale morphism \(q\) is an open immersion,
the locus on which \(f'\) and \(g'\) agree is represented by an open
immersion
\[
    Y_0
    \hookrightarrow
    Y.
\]
The prescribed trivializations and the Cartesian condition over \(Z\)
imply that \(f'\) and \(g'\) agree after restriction to \(Y_Z\).
Hence
\[
    Y_Z
    \hookrightarrow
    Y
\]
factors through \(Y_0\).

It follows that
\[
    \{Y_0\to Y,\;Y_U\to Y\}
\]
is a Zariski covering family.  The two sections agree over \(Y_0\),
while over \(Y_U\) the relevant trivialized-map space is contractible.
Thus the diagonal of
\[
    L_{\mathrm{Nis},S}(\varphi)
\]
is an effective epimorphism.

We have shown that \(L_{\mathrm{Nis},S}(\varphi)\) is \(0\)-truncated,
that it is an effective epimorphism, and that its diagonal is an
effective epimorphism.  Since the diagonal of a \(0\)-truncated
morphism is \((-1)\)-truncated, its being an effective epimorphism
forces it to be an equivalence.  Therefore
\[
    L_{\mathrm{Nis},S}(\varphi)
\]
is a monomorphism.  Being both a monomorphism and an effective
epimorphism, it is an equivalence.

Hence
\[
    h_S(X',t')
    \longrightarrow
    h_S(X,t)
\]
is a Nisnevich-local equivalence and consequently a motivic
equivalence.
\end{proof}

\subsection{Vector-bundle contraction}

\begin{proposition}[Contractibility of the vector-bundle model]
\label{prop:vector-bundle-trivialized-maps-contractible}
Let
\[
    \pi
    \colon
    \mathbb V_S(E)
    \longrightarrow
    S
\]
be the relative linear space associated with a finite locally free
quasi-coherent module \(E\), and let
\[
    s
    \colon
    S
    \longrightarrow
    \mathbb V_S(E)
\]
be its zero section.  Write
\[
    s_Z
    \colon
    Z
    \longrightarrow
    \mathbb V_S(E)
\]
for its restriction.  Then
\[
    L_{\mathrm{mot},S}
    h_S\bigl(
        \mathbb V_S(E),
        s_Z
    \bigr)
    \simeq
    *_S.
\]
\end{proposition}

\begin{proof}
Scalar multiplication defines a morphism
\[
    m
    \colon
    \mathbb A^1_S
    \times_S
    \mathbb V_S(E)
    \longrightarrow
    \mathbb V_S(E).
\]
If
\[
    (f,\eta)
    \in
    h_S(\mathbb V_S(E),s_Z)(Y),
\]
then the composite
\[
    \mathbb A^1_Y
    \xrightarrow{\id\times f}
    \mathbb A^1_S\times_S\mathbb V_S(E)
    \xrightarrow{m}
    \mathbb V_S(E)
\]
remains equal to the zero section after restriction to
\(\mathbb A^1_{Y_Z}\).  Hence scalar multiplication induces an
\(\mathbb A^1\)-homotopy on
\[
    h_S(\mathbb V_S(E),s_Z).
\]
At the unit section \(1\), this homotopy is the identity.  At the zero
section \(0\), it is the constant map through the global zero section
\(s\).  After imposing \(\mathbb A^1\)-invariance, the identity is
therefore homotopic to a constant map.  The global zero section supplies
a point, so the motivic localization is contractible.
\end{proof}

\subsection{The partial-section theorem}

\begin{lemma}[Cartesian refinement of a spectral tubular neighbourhood]
\label{lem:cartesian-tubular-neighbourhood}
Let
\[
    p
    \colon
    X
    \longrightarrow
    S
\]
be a smooth morphism of finite presentation, and let
\[
    s
    \colon
    S
    \longrightarrow
    X
\]
be a section.  Put
\[
    N
    :=
    N^*_{S/X}.
\]
There exists a Zariski open cover
\[
    \{S_\alpha\to S\}_\alpha
\]
such that, after writing
\[
    X_\alpha
    :=
    X\times_SS_\alpha,
\]
there are an open neighbourhood
\[
    X_\alpha^\circ
    \subseteq
    X_\alpha
\]
of \(s(S_\alpha)\) and an étale morphism
\[
    \varphi_\alpha
    \colon
    X_\alpha^\circ
    \longrightarrow
    \mathbb V_{S_\alpha}
    \bigl(
        N|_{S_\alpha}
    \bigr)
\]
for which the square
\[
\begin{tikzcd}[column sep=large, row sep=large]
    S_\alpha
        \arrow[r, "s"]
        \arrow[d, "\id_{S_\alpha}"']
    &
    X_\alpha^\circ
        \arrow[d, "\varphi_\alpha"]
    \\
    S_\alpha
        \arrow[r, "0"']
    &
    \mathbb V_{S_\alpha}
    \bigl(
        N|_{S_\alpha}
    \bigr)
\end{tikzcd}
\]
is Cartesian.
\end{lemma}

\begin{proof}
Fix a point
\[
    x\in|S|.
\]
Choose an affine open neighbourhood
\[
    S_0\subseteq S
\]
of \(x\).  Choose an affine open neighbourhood
\[
    W\subseteq X\times_SS_0
\]
of the point \(s(x)\).

The inverse image
\[
    s^{-1}(W)
    \subseteq
    S_0
\]
is an open neighbourhood of \(x\).  After replacing \(S_0\) by a
smaller affine open neighbourhood contained in \(s^{-1}(W)\), the
section \(s|_{S_0}\) factors through
\[
    W
    \subseteq
    X\times_SS_0.
\]
We continue to denote this affine neighbourhood by \(W\).  Since both
\(W\) and \(S_0\) are affine, the morphism
\[
    W
    \longrightarrow
    S_0
\]
is affine.  It is smooth, being the restriction of \(p\).

The conormal module of the section computed in \(W\) agrees with the
restriction of the conormal module computed in \(X\):
\[
    N^*_{S_0/W}
    \simeq
    N|_{S_0},
\]
because \(W\hookrightarrow X\times_SS_0\) is an open immersion
containing the image of the section.

Apply
\Ref{thm:spectral-tubular-neighbourhood}
to the smooth affine morphism
\[
    W
    \longrightarrow
    S_0
\]
and its section.  After shrinking \(S_0\) further if necessary, there
are an open neighbourhood
\[
    W^\circ
    \subseteq
    W
\]
of \(s(S_0)\) and an étale morphism
\[
    \varphi
    \colon
    W^\circ
    \longrightarrow
    \mathbb V_{S_0}
    \bigl(
        N|_{S_0}
    \bigr)
\]
which carries \(s\) to the zero section.

Form the zero fibre
\[
    P
    :=
    W^\circ
    \times_{
        \mathbb V_{S_0}(N|_{S_0})
    }
    S_0.
\]
The morphism
\[
    P
    \longrightarrow
    S_0
\]
is étale, and the section \(s\) induces a section
\[
    \sigma
    \colon
    S_0
    \longrightarrow
    P.
\]
Since \(P\to S_0\) is étale, the section \(\sigma\) is an open
immersion.

It is also a closed immersion.  Indeed, the morphism
\[
    W^\circ
    \longrightarrow
    S_0
\]
is separated, because it is the restriction of the affine morphism
\(W\to S_0\).  Hence its section
\[
    s
    \colon
    S_0
    \longrightarrow
    W^\circ
\]
is a closed immersion.  The square
\[
\begin{tikzcd}[column sep=large, row sep=large]
    S_0
        \arrow[r, "\sigma"]
        \arrow[d, "\id_{S_0}"']
    &
    P
        \arrow[d]
    \\
    S_0
        \arrow[r, "s"']
    &
    W^\circ
\end{tikzcd}
\]
is Cartesian, since the zero section of
\(\mathbb V_{S_0}(N|_{S_0})\) is a monomorphism.  Thus \(\sigma\) is
closed.

Consequently, \(\sigma\) is open and closed, and there is a
decomposition
\[
    P
    \simeq
    S_0
    \coprod
    P',
\]
where the first summand is the image of \(\sigma\).

The zero section of
\[
    \mathbb V_{S_0}(N|_{S_0})
\]
is a closed immersion, so
\[
    P
    \hookrightarrow
    W^\circ
\]
is a closed immersion.  Hence the composite
\[
    P'
    \hookrightarrow
    W^\circ
\]
is closed.  Put
\[
    W^{\circ\circ}
    :=
    W^\circ\setminus P'.
\]
Then \(W^{\circ\circ}\) is an open neighbourhood of \(s(S_0)\), the
restriction
\[
    W^{\circ\circ}
    \longrightarrow
    \mathbb V_{S_0}
    \bigl(
        N|_{S_0}
    \bigr)
\]
remains étale, and its zero fibre is precisely \(S_0\):
\[
    W^{\circ\circ}
    \times_{
        \mathbb V_{S_0}(N|_{S_0})
    }
    S_0
    \simeq
    S_0.
\]

Repeating this construction around every point of \(S\) gives the
required Zariski open cover and the asserted Cartesian étale
neighbourhoods.
\end{proof}


\begin{lemma}[Nisnevich extension of a partial section]
\label{lem:nisnevich-extension-partial-section}
Let
\[
    i
    \colon
    Z
    \hookrightarrow
    S
\]
be a closed immersion, let
\[
    p
    \colon
    X
    \longrightarrow
    S
\]
be smooth and of finite presentation, and let
\[
    t
    \colon
    Z
    \longrightarrow
    X
\]
be an \(S\)-morphism.

There exists a finite family of étale morphisms
\[
    q_a
    \colon
    Y_a
    \longrightarrow
    S,
    \qquad
    a\in A,
\]
and \(S\)-morphisms
\[
    r_a
    \colon
    Y_a
    \longrightarrow
    X
\]
such that:
\begin{enumerate}[label=(\roman*)]
    \item the family
    \[
        \{(Y_a)_Z\to Z\}_{a\in A}
    \]
    is a Nisnevich covering family;

    \item for every \(a\), the restriction of \(r_a\) to
    \((Y_a)_Z\) agrees with the composite
    \[
        (Y_a)_Z
        \longrightarrow
        Z
        \xrightarrow{t}
        X.
    \]
\end{enumerate}
Consequently,
\[
    \{U\to S,\;Y_a\to S\}_{a\in A},
    \qquad
    U=S\setminus Z,
\]
is a spectral Nisnevich covering family of \(S\).
\end{lemma}

\begin{proof}
Fix a point
\[
    z\in|Z|.
\]
After replacing \(S\) by an affine open neighbourhood of the image of
\(z\), and replacing \(X\) by an affine open neighbourhood of \(t(z)\),
we may suppose that \(S\) and \(X\) are affine and that
\[
    X
    \longrightarrow
    S
\]
is smooth affine.  Shrinking the chosen neighbourhood of \(S\)
further, we may also suppose that the restricted partial section
factors through this affine open of \(X\).

Apply
\Ref{lem:cartesian-tubular-neighbourhood}
to the smooth morphism
\[
    X_Z
    \longrightarrow
    Z
\]
and its section \(t\).  After a further Zariski localization around
\(z\), there are a finite locally free module \(N\) on \(Z\), an open
neighbourhood
\[
    X_Z^\circ
    \subseteq
    X_Z
\]
of \(t(Z)\), and an étale morphism
\[
    f
    \colon
    X_Z^\circ
    \longrightarrow
    \mathbb V_Z(N)
\]
which is Cartesian over the zero section.

After trivializing \(N\) Zariski-locally and shrinking once more, write
\[
    \mathbb V_Z(N)
    \simeq
    \mathbb A_Z^n.
\]
Replace \(X\) by an affine open neighbourhood whose closed base change
is contained in \(X_Z^\circ\).  The morphism
\[
    f
    \colon
    X_Z
    \longrightarrow
    \mathbb A_Z^n
\]
is determined by finitely many points of the connective coordinate
algebra of \(X_Z\).  Since
\[
    X_Z
    \hookrightarrow
    X
\]
is a closed immersion, the induced map on degree-zero homotopy groups
of coordinate algebras is surjective.  Choose lifts of the
corresponding components, together with paths identifying their
restrictions with the original points.  They define an \(S\)-morphism
\[
    g
    \colon
    X
    \longrightarrow
    \mathbb A_S^n
\]
whose base change to \(Z\) is identified with \(f\).

Define
\[
    Y_0
    :=
    X\times_{\mathbb A_S^n}S,
\]
where
\[
    S
    \longrightarrow
    \mathbb A_S^n
\]
is the zero section.  Base change to \(Z\) gives
\[
\begin{aligned}
    (Y_0)_Z
    &\simeq
    X_Z\times_{\mathbb A_Z^n}Z
    \\
    &\simeq
    Z.
\end{aligned}
\]
The resulting equivalence is induced by the partial section \(t\).

The morphism
\[
    Y_0
    \longrightarrow
    S
\]
is locally of finite presentation.  Indeed, locally it is obtained by
base change from the morphism
\[
    X
    \longrightarrow
    \mathbb A_S^n,
\]
which is locally of finite presentation because both its source and
target are of finite presentation over \(S\).

Cotangent base change gives
\[
    L_{Y_0/S}|_Z
    \simeq
    L_{X_Z/\mathbb A_Z^n}|_Z
    \simeq
    0,
\]
since \(f\) is étale.

Because
\[
    Y_0
    \longrightarrow
    S
\]
is locally of finite presentation, its relative cotangent complex
\[
    L_{Y_0/S}
\]
is perfect.  The support of a perfect quasi-coherent module is closed:
affine-locally, compactness implies that vanishing after localization at
a point already holds on a principal neighbourhood.  Consequently, the
vanishing locus
\[
    E
    :=
    |Y_0|
    \setminus
    \operatorname{Supp}(L_{Y_0/S})
\]
is represented by an open spectral subscheme
\[
    E
    \subseteq
    Y_0.
\]
The preceding cotangent calculation shows that
\[
    (Y_0)_Z
    \subseteq
    E.
\]

On \(E\), the restricted morphism
\[
    E
    \longrightarrow
    S
\]
is locally of finite presentation and has vanishing relative cotangent
complex.  It is therefore étale.

Choose an affine open neighbourhood
\[
    Y
    \subseteq
    E
\]
of the point corresponding to \(z\) under
\[
    (Y_0)_Z
    \simeq
    Z.
\]
In particular, \(Y\hookrightarrow Y_0\) is a quasi-compact open
immersion and hence is of finite presentation.  The composite
\[
    q
    \colon
    Y
    \longrightarrow
    Y_0
    \longrightarrow
    S
\]
is therefore étale and of finite presentation.

Under the equivalence
\[
    (Y_0)_Z
    \simeq
    Z,
\]
the closed base change
\[
    Y_Z
    :=
    Y\times_SZ
\]
is an open neighbourhood of \(z\) in \(Z\).  In particular,
\[
    Y_Z
    \longrightarrow
    Z
\]
is a Nisnevich neighbourhood of \(z\).

The projection
\[
    r
    \colon
    Y
    \longrightarrow
    Y_0
    \longrightarrow
    X
\]
restricts over \(Y_Z\) to the composite
\[
    Y_Z
    \longrightarrow
    Z
    \xrightarrow{t}
    X.
\]

Repeating the construction around every point of \(Z\) gives a
Nisnevich covering family
\[
    \{(Y_a)_Z\to Z\}_{a\in A}.
\]
Since \(Z\) is quasi-compact, a finite subfamily suffices.

Adjoining the open complement
\[
    U
    \hookrightarrow
    S
\]
gives a spectral Nisnevich covering family
\[
    \{U\to S,\;Y_a\to S\}_{a\in A},
\]
because the étale family is Nisnevich-surjective over the closed
complement.

See
\cite{KhanBraveNewMotivicI}.
\end{proof}


\begin{theorem}[Motivic contractibility of partial-section spaces]
\label{thm:partial-section-motivic-contractibility}
Let
\[
    X
    \longrightarrow
    S
\]
be smooth and of finite presentation, and let
\[
    t
    \colon
    Z
    \longrightarrow
    X
\]
be an \(S\)-morphism.  Then
\[
    L_{\mathrm{mot},S}h_S(X,t)
    \simeq
    *_S.
\]
\end{theorem}

\begin{proof}
Apply
\Ref{lem:nisnevich-extension-partial-section}
to obtain a finite Nisnevich covering family
\[
    \{U\to S,\;q_a:Y_a\to S\}_{a\in A}
\]
such that each \(q_a\) admits an \(S\)-morphism
\[
    r_a
    \colon
    Y_a
    \longrightarrow
    X
\]
whose restriction to
\[
    (Y_a)_Z
    =
    Y_a\times_SZ
\]
agrees with the pullback of \(t\).

By
\Ref{prop:nisnevich-local-conservativity},
it is enough to prove contractibility after pullback to \(U\) and to
each \(Y_a\).

Over \(U\), the closed base change is empty.  By the definition of the
trivialized-map presheaf,
\[
    h_U(X_U,t_U)
    \simeq
    *_U.
\]

Fix \(a\), and write
\[
    Y:=Y_a,
    \qquad
    r:=r_a.
\]
Smooth base change gives
\[
    q_a^*L_{\mathrm{mot},S}h_S(X,t)
    \simeq
    L_{\mathrm{mot},Y}h_Y(X_Y,t_Y),
\]
where
\[
    X_Y
    :=
    X\times_SY.
\]

The morphism
\[
    r
    \colon
    Y
    \longrightarrow
    X
\]
determines a section
\[
    \sigma
    \colon
    Y
    \longrightarrow
    X_Y,
    \qquad
    y
    \longmapsto
    \bigl(
        r(y),y
    \bigr),
\]
whose restriction to
\[
    Y_Z
    :=
    Y\times_SZ
\]
is \(t_Y\).

The morphism
\[
    X_Y
    \longrightarrow
    Y
\]
is smooth and of finite presentation.  Apply
\Ref{lem:cartesian-tubular-neighbourhood}
to this morphism and the section \(\sigma\).  There is a Zariski cover
\[
    \{Y_\beta\to Y\}_\beta
\]
such that, after putting
\[
    X_\beta
    :=
    X_Y\times_YY_\beta,
\]
there are an open neighbourhood
\[
    X_\beta^\circ
    \subseteq
    X_\beta
\]
of \(\sigma(Y_\beta)\), a finite locally free module \(N_\beta\) on
\(Y_\beta\), and an étale morphism
\[
    \varphi_\beta
    \colon
    X_\beta^\circ
    \longrightarrow
    \mathbb V_{Y_\beta}(N_\beta)
\]
such that
\[
\begin{tikzcd}[column sep=large, row sep=large]
    Y_\beta
        \arrow[r, "\sigma"]
        \arrow[d, "\id_{Y_\beta}"']
    &
    X_\beta^\circ
        \arrow[d, "\varphi_\beta"]
    \\
    Y_\beta
        \arrow[r, "0"']
    &
    \mathbb V_{Y_\beta}(N_\beta)
\end{tikzcd}
\]
is Cartesian.

For each \(\beta\), smooth base change for trivialized-map presheaves
gives
\[
    L_{\mathrm{mot},Y_\beta}
    h_{Y_\beta}
    \bigl(
        X_\beta,
        t_{Y_\beta}
    \bigr)
\]
as the pullback of
\[
    L_{\mathrm{mot},Y}
    h_Y(X_Y,t_Y).
\]

The open immersion
\[
    X_\beta^\circ
    \hookrightarrow
    X_\beta
\]
is Cartesian over the section \(\sigma\), and the étale morphism
\[
    \varphi_\beta
    \colon
    X_\beta^\circ
    \longrightarrow
    \mathbb V_{Y_\beta}(N_\beta)
\]
is Cartesian over the zero section.  Therefore
\Ref{prop:etale-invariance-trivialized-maps}
gives motivic equivalences
\[
\begin{aligned}
    h_{Y_\beta}
    \bigl(
        X_\beta,
        t_{Y_\beta}
    \bigr)
    &\simeq_{\mathrm{mot}}
    h_{Y_\beta}
    \bigl(
        X_\beta^\circ,
        \sigma|_{(Y_\beta)_Z}
    \bigr)
    \\
    &\simeq_{\mathrm{mot}}
    h_{Y_\beta}
    \bigl(
        \mathbb V_{Y_\beta}(N_\beta),
        0_{(Y_\beta)_Z}
    \bigr).
\end{aligned}
\]
By
\Ref{prop:vector-bundle-trivialized-maps-contractible},
\[
    L_{\mathrm{mot},Y_\beta}
    h_{Y_\beta}
    \bigl(
        \mathbb V_{Y_\beta}(N_\beta),
        0_{(Y_\beta)_Z}
    \bigr)
    \simeq
    *_{Y_\beta}.
\]
Thus the pullback of
\[
    L_{\mathrm{mot},Y}h_Y(X_Y,t_Y)
\]
to every \(Y_\beta\) is terminal.

The family
\[
    \{Y_\beta\to Y\}_\beta
\]
is a Nisnevich covering family.  By
\Ref{prop:nisnevich-local-conservativity},
\[
    L_{\mathrm{mot},Y}h_Y(X_Y,t_Y)
    \simeq
    *_Y.
\]

The original Nisnevich covering family of \(S\), together with another
application of
\Ref{prop:nisnevich-local-conservativity},
now gives
\[
    L_{\mathrm{mot},S}h_S(X,t)
    \simeq
    *_S.
\]
\end{proof}

\subsection{Localization on standard smooth generators}

\begin{lemma}[Presheaf projection and smooth--closed exchange]
\label{lem:presheaf-projection-smooth-closed-exchange}
Let
\[
    p
    \colon
    Y
    \longrightarrow
    S
\]
be smooth and of finite presentation.

For presheaves
\[
    A
    \in
    \operatorname{PSh}(Y),
    \qquad
    G
    \in
    \operatorname{PSh}(S),
\]
there is a canonical presheaf projection equivalence
\[
    p^{\operatorname{PSh}}_\sharp
    \bigl(
        A\times p_{\operatorname{PSh}}^*G
    \bigr)
    \simeq
    p^{\operatorname{PSh}}_\sharp A
    \times
    G.
\]
The corresponding relative form is the following.  If
\[
    G
    \longrightarrow
    H
\]
is a morphism in \(\operatorname{PSh}(S)\), and
\[
    A
    \longrightarrow
    p^*H
\]
is a morphism in \(\operatorname{PSh}(Y)\), then
\[
    p^{\operatorname{PSh}}_\sharp
    \bigl(
        A\times_{p^*H}p^*G
    \bigr)
    \simeq
    p^{\operatorname{PSh}}_\sharp A
    \times_H
    G.
\]

For the Cartesian square
\[
\begin{tikzcd}[column sep=large, row sep=large]
    Y_Z
        \arrow[r, "q"]
        \arrow[d, "k"']
    &
    Z
        \arrow[d, "i"]
    \\
    Y
        \arrow[r, "p"']
    &
    S,
\end{tikzcd}
\]
there is also a canonical presheaf-level exchange equivalence
\[
    p_{\operatorname{PSh}}^*
    i_{*,\operatorname{PSh}}
    \simeq
    k_{*,\operatorname{PSh}}
    q_{\operatorname{PSh}}^*.
\]
\end{lemma}

\begin{proof}
The two sides of the first comparison preserve colimits separately in
\(A\) and \(G\).  It is therefore enough to take
\[
    A=h_Y(T),
    \qquad
    G=h_S(V),
\]
with \(T\in\mathcal C_Y\) and \(V\in\mathcal C_S\).  Then
\[
\begin{aligned}
    p^{\operatorname{PSh}}_\sharp
    \bigl(
        h_Y(T)\times p^*h_S(V)
    \bigr)
    &\simeq
    p^{\operatorname{PSh}}_\sharp
    h_Y(T\times_SV)
    \\
    &\simeq
    h_S(T\times_SV),
\end{aligned}
\]
while
\[
\begin{aligned}
    p^{\operatorname{PSh}}_\sharp h_Y(T)
    \times h_S(V)
    &\simeq
    h_S(T)\times h_S(V)
    \\
    &\simeq
    h_S(T\times_SV).
\end{aligned}
\]
The comparison is the identity under these identifications.  Applying
the same argument in the presheaf slice categories gives the relative
form.

For the exchange equivalence, let
\[
    F
    \in
    \operatorname{PSh}(Z)
\]
and let \(T\in\mathcal C_Y\).  Since smooth presheaf inverse image is
restriction along composition,
\[
    \bigl(
        p^*i_*F
    \bigr)(T)
    \simeq
    F(T\times_SZ).
\]
On the other hand,
\[
\begin{aligned}
    \bigl(
        k_*q^*F
    \bigr)(T)
    &\simeq
    \bigl(
        q^*F
    \bigr)(T\times_YY_Z)
    \\
    &\simeq
    F(T\times_YY_Z)
    \\
    &\simeq
    F(T\times_SZ).
\end{aligned}
\]
These objectwise equivalences are natural and give the asserted
exchange equivalence.
\end{proof}


\begin{proposition}[Localization for a standard smooth motive]
\label{prop:localization-standard-smooth-motive}
Let
\[
    X
    \longrightarrow
    S
\]
be standard smooth.  Then the canonical morphism
\[
    L_{\mathrm{mot},S}h_S^Z(X)
    \longrightarrow
    i_*\mathrm M_Z(X_Z)
\]
is an equivalence.
\end{proposition}

\begin{proof}
The morphism in the statement is the composite
\[
\begin{aligned}
    L_{\mathrm{mot},S}h_S^Z(X)
    &\xrightarrow{
        L_{\mathrm{mot},S}(\rho_X)
    }
    L_{\mathrm{mot},S}
    i_{*,\operatorname{PSh}}h_Z(X_Z)
    \\
    &\xrightarrow{\simeq}
    i_*\mathrm M_Z(X_Z),
\end{aligned}
\]
where the second morphism is the equivalence of
\Ref{cor:closed-restriction-motivic-localization}.
It is therefore enough to prove that
\[
    \rho_X
    \colon
    h_S^Z(X)
    \longrightarrow
    i_{*,\operatorname{PSh}}h_Z(X_Z)
\]
is a motivic equivalence.

By the co-Yoneda presentation,
\[
    i_{*,\operatorname{PSh}}h_Z(X_Z)
    \simeq
    \operatorname*{colim}_{(Y,\tau)}
    h_S(Y),
\]
where the indexing category consists of pairs
\[
    Y
    \in
    \mathcal C_S,
    \qquad
    \tau
    \colon
    Y_Z
    \longrightarrow
    X_Z.
\]
Colimits in a presheaf category are universal.  Hence
\[
    h_S^Z(X)
    \simeq
    \operatorname*{colim}_{(Y,\tau)}
    \left(
        h_S^Z(X)
        \times_{
            i_{*,\operatorname{PSh}}h_Z(X_Z)
        }
        h_S(Y)
    \right).
\]
Under these two colimit presentations, \(\rho_X\) is the colimit of the
projection morphisms
\[
    h_S^Z(X)
    \times_{
        i_{*,\operatorname{PSh}}h_Z(X_Z)
    }
    h_S(Y)
    \longrightarrow
    h_S(Y).
\]
It is therefore enough to prove that every such projection is a
motivic equivalence.

Let
\[
    p
    \colon
    Y
    \longrightarrow
    S
\]
be the structure morphism, and form
\[
\begin{tikzcd}[column sep=large, row sep=large]
    Y_Z
        \arrow[r, "q"]
        \arrow[d, "k"']
    &
    Z
        \arrow[d, "i"]
    \\
    Y
        \arrow[r, "p"']
    &
    S.
\end{tikzcd}
\]
The map \(\tau\) determines a partial section
\[
    t_\tau
    \colon
    Y_Z
    \longrightarrow
    X_Y
    :=
    X\times_SY,
    \qquad
    y
    \longmapsto
    \bigl(
        \tau(y),y
    \bigr).
\]

By
\Ref{lem:presheaf-projection-smooth-closed-exchange},
the projection under consideration is canonically identified with the
image under
\[
    p^{\operatorname{PSh}}_\sharp
\]
of the structure morphism
\[
    h_Y(X_Y,t_\tau)
    \longrightarrow
    *_Y.
\]
Indeed, smooth base change identifies
\[
    p^*h_S^Z(X)
    \simeq
    h_Y^{Y_Z}(X_Y),
\]
while presheaf smooth--closed exchange identifies
\[
    p^*
    i_{*,\operatorname{PSh}}h_Z(X_Z)
    \simeq
    k_{*,\operatorname{PSh}}
    h_{Y_Z}(X_Z\times_ZY_Z).
\]
Under these identifications, the morphism supplied by \(\tau\) is the
point defining
\[
    h_Y(X_Y,t_\tau).
\]
The relative presheaf projection formula therefore gives
\[
\begin{aligned}
    &
    h_S^Z(X)
    \times_{
        i_{*,\operatorname{PSh}}h_Z(X_Z)
    }
    h_S(Y)
    \\
    &\hspace{3em}\simeq
    p^{\operatorname{PSh}}_\sharp
    h_Y(X_Y,t_\tau).
\end{aligned}
\]

The morphism
\[
    X_Y
    \longrightarrow
    Y
\]
is smooth and of finite presentation.  Therefore
\Ref{thm:partial-section-motivic-contractibility}
gives
\[
    L_{\mathrm{mot},Y}h_Y(X_Y,t_\tau)
    \simeq
    *_Y.
\]
Smooth direct image descends through motivic localization and preserves
equivalences.  Hence
\[
    L_{\mathrm{mot},S}
    p^{\operatorname{PSh}}_\sharp
    h_Y(X_Y,t_\tau)
    \longrightarrow
    L_{\mathrm{mot},S}
    p^{\operatorname{PSh}}_\sharp(*_Y)
\]
is an equivalence.  Since
\[
    p^{\operatorname{PSh}}_\sharp(*_Y)
    \simeq
    h_S(Y),
\]
the projection
\[
    h_S^Z(X)
    \times_{
        i_{*,\operatorname{PSh}}h_Z(X_Z)
    }
    h_S(Y)
    \longrightarrow
    h_S(Y)
\]
is a motivic equivalence.

Motivic equivalences are stable under colimits.  Taking the colimit over
all pairs \((Y,\tau)\) therefore shows that
\[
    \rho_X
    \colon
    h_S^Z(X)
    \longrightarrow
    i_{*,\operatorname{PSh}}h_Z(X_Z)
\]
is a motivic equivalence.  Consequently,
\[
    L_{\mathrm{mot},S}(\rho_X)
\]
is an equivalence.  Composing with
\[
    L_{\mathrm{mot},S}
    i_{*,\operatorname{PSh}}h_Z(X_Z)
    \simeq
    i_*\mathrm M_Z(X_Z)
\]
from
\Ref{cor:closed-restriction-motivic-localization}
gives the asserted equivalence
\[
    L_{\mathrm{mot},S}h_S^Z(X)
    \xrightarrow{\simeq}
    i_*\mathrm M_Z(X_Z).
\]
\end{proof}

\subsection{Proof of localization}

\begin{proof}[Proof of \Ref{thm:spectral-morel-voevodsky-localization}]
Let
\[
    \Lambda(F)
    :=
    F
    \amalg_{j_\sharp j^*F}
    j_\sharp(*_U).
\]
The localization transformation is the natural morphism
\[
    \lambda_F
    \colon
    \Lambda(F)
    \longrightarrow
    i_*i^*F.
\]

The functors \(j^*\), \(j_\sharp\), and \(i^*\) preserve all colimits.
The functor \(i_*\) preserves weakly contractible colimits by
\Ref{thm:closed-direct-image-weakly-contractible-colimits}.  Every
sifted indexing category is weakly contractible.  Moreover, for a
sifted indexing category \(K\),
\[
    \operatorname*{colim}_{k\in K}j_\sharp(*_U)
    \simeq
    j_\sharp(*_U).
\]
It follows by universality of colimits that both
\[
    \Lambda
    \quad\text{and}\quad
    i_*i^*
\]
preserve sifted colimits.

By \Ref{prop:generation-by-standard-smooth-motives}, it is enough to
prove that \(\lambda_F\) is an equivalence for
\[
    F
    =
    \mathrm M_S(X)
\]
with \(X\) standard smooth.  For such \(X\),
\Ref{prop:localization-square-smooth-motive} identifies \(\lambda_F\)
with
\[
    L_{\mathrm{mot},S}h_S^Z(X)
    \longrightarrow
    i_*\mathrm M_Z(X_Z),
\]
which is an equivalence by
\Ref{prop:localization-standard-smooth-motive}.  Hence the localization
square is cocartesian for every \(F\in\mathbf H(S)\).
\end{proof}

\section{Closed recollement and exchange}
\label{sec:closed-recollement-exchange}

\subsection{Pointed localization sequences}

\begin{theorem}[Pointed localization cofibre sequence]
\label{thm:pointed-localization-cofiber-sequence}
For every
\[
    F
    \in
    \mathbf H_\bullet(S),
\]
there is a natural cofibre sequence
\[
    j_\sharp j^*F
    \longrightarrow
    F
    \longrightarrow
    i_*i^*F.
\]
\end{theorem}

\begin{proof}
Let \(U(F)\) denote the underlying unpointed motivic space.  The pointed
smooth direct image is defined by the pushout
\[
    j_{\sharp,\bullet}(j^*F)
    =
    j_\sharp U(j^*F)
    \amalg_{j_\sharp(*_U)}
    *_S.
\]
Paste this defining pushout with the unpointed localization square of
\(U(F)\).  The result is a cocartesian square in
\(\mathbf H_\bullet(S)\):
\[
\begin{tikzcd}[column sep=large, row sep=large]
    j_{\sharp,\bullet}j^*F
        \arrow[r]
        \arrow[d]
    &
    F
        \arrow[d]
    \\
    0
        \arrow[r]
    &
    i_*i^*F.
\end{tikzcd}
\]
This is exactly the asserted cofibre sequence.  We suppress the
subscript \(\bullet\) on \(j_\sharp\) in the statement.
\end{proof}

\subsection{Closed support and full faithfulness}

\begin{lemma}[Closed-liftable Nisnevich basis]
\label{lem:closed-liftable-nisnevich-basis}
Let
\[
    i
    \colon
    Z
    \hookrightarrow
    S
\]
be a closed immersion.  Let
\[
    \mathcal B_i
    \subseteq
    \operatorname{Sm}^{\mathrm{fp}}_{/Z}
\]
be the full subcategory spanned by those \(Y\to Z\) for which there
exists
\[
    \widetilde Y
    \in
    \operatorname{Sm}^{\mathrm{fp}}_{/S}
\]
and an equivalence
\[
    \widetilde Y\times_SZ
    \simeq
    Y.
\]

Then \(\mathcal B_i\) is a basis for the spectral Nisnevich topology on
\(\operatorname{Sm}^{\mathrm{fp}}_{/Z}\).  In particular, a morphism of
spectral Nisnevich sheaves is an equivalence if and only if it is an
equivalence when evaluated on every object of \(\mathcal B_i\).
\end{lemma}

\begin{proof}
Let
\[
    Y
    \longrightarrow
    Z
\]
be smooth and of finite presentation.  By
\Ref{prop:local-lifting-smooth-test-schemes},
every point of \(Y\) has a Nisnevich neighbourhood
\[
    Y_\alpha
    \longrightarrow
    Y
\]
which belongs to \(\mathcal B_i\).  Since \(Y\) is quasi-compact, a
finite subfamily covers \(Y\).

Let
\[
    B_1
    \longrightarrow
    Y,
    \qquad
    B_2
    \longrightarrow
    Y
\]
be morphisms with \(B_1,B_2\in\mathcal B_i\).  Their fibre product
\[
    B_1\times_YB_2
\]
is smooth and of finite presentation over \(Z\).  Applying
\Ref{prop:local-lifting-smooth-test-schemes} to this fibre product gives
a finite Nisnevich covering by objects of \(\mathcal B_i\).  The same
argument applies after arbitrary pullback of a covering family by
objects of \(\mathcal B_i\).  Thus every object admits a covering by
objects of \(\mathcal B_i\), and pullbacks of basis coverings admit
refinements by basis coverings.  This proves the basis axioms.

Let
\[
    v
    \colon
    \mathcal B_i
    \hookrightarrow
    \operatorname{Sm}^{\mathrm{fp}}_{/Z}
\]
be the inclusion.  By the comparison lemma for a basis, restriction
induces an equivalence
\[
    v^*
    \colon
    \operatorname{Sh}_{\mathrm{Nis}}
    \bigl(
        \operatorname{Sm}^{\mathrm{fp}}_{/Z}
    \bigr)
    \xrightarrow{\simeq}
    \operatorname{Sh}_{\mathrm{Nis}}(\mathcal B_i).
\]
Consequently, a morphism of spectral Nisnevich sheaves is an
equivalence if and only if its restriction to \(\mathcal B_i\) is an
equivalence, equivalently if and only if it is an equivalence when
evaluated on every object of \(\mathcal B_i\).
\end{proof}

\begin{lemma}[Conservativity of closed direct image]
\label{lem:closed-direct-image-conservative}
The functor
\[
    i_*
    \colon
    \mathbf H(Z)
    \longrightarrow
    \mathbf H(S)
\]
is conservative.  The same is true on pointed motivic spaces.
\end{lemma}

\begin{proof}
Let
\[
    \varphi
    \colon
    G
    \longrightarrow
    H
\]
be a morphism in \(\mathbf H(Z)\) such that
\[
    i_*\varphi
\]
is an equivalence.

Let
\[
    Y
    \in
    \mathcal B_i
\]
be an object of the closed-liftable Nisnevich basis of
\Ref{lem:closed-liftable-nisnevich-basis}.  Choose
\[
    \widetilde Y
    \in
    \operatorname{Sm}^{\mathrm{fp}}_{/S}
\]
and an equivalence
\[
    \widetilde Y\times_SZ
    \simeq
    Y.
\]
The evaluation formula for direct image gives
\[
\begin{aligned}
    G(Y)
    &\simeq
    (i_*G)(\widetilde Y),
    \\
    H(Y)
    &\simeq
    (i_*H)(\widetilde Y).
\end{aligned}
\]
Under these identifications,
\[
    \varphi(Y)
\]
is the evaluation of \(i_*\varphi\) on \(\widetilde Y\).  It is
therefore an equivalence.

Thus \(\varphi\) is an equivalence on every object of the basis
\(\mathcal B_i\).  By
\Ref{lem:closed-liftable-nisnevich-basis},
it is an equivalence of Nisnevich sheaves and hence an equivalence in
\(\mathbf H(Z)\).

The pointed forgetful functor reflects equivalences, so the pointed
assertion follows.
\end{proof}

\begin{theorem}[Closed support theorem]
\label{thm:closed-support-kashiwara}
The closed direct image
\[
    i_*
    \colon
    \mathbf H(Z)
    \longrightarrow
    \mathbf H(S)
\]
is fully faithful.  Its essential image is
\[
    \mathbf H_Z(S).
\]
Likewise,
\[
    i_*
    \colon
    \mathbf H_\bullet(Z)
    \longrightarrow
    \mathbf H_\bullet(S)
\]
is fully faithful with essential image
\[
    \mathbf H_{\bullet,Z}(S).
\]
\end{theorem}

\begin{proof}
Let \(G\in\mathbf H(Z)\).  Apply localization to \(i_*G\).  By
\Ref{prop:open-closed-orthogonality},
\[
    j^*i_*G
    \simeq
    *_U.
\]
The two terms in the left column of the localization square are
therefore equivalent, and localization gives an equivalence
\[
    \eta_{i_*G}
    \colon
    i_*G
    \xrightarrow{\simeq}
    i_*i^*i_*G.
\]
Let
\[
    \varepsilon_G
    \colon
    i^*i_*G
    \longrightarrow
    G
\]
be the counit of \(i^*\dashv i_*\).  The triangle identity identifies
\[
    i_*(\varepsilon_G)
\]
with the inverse of \(\eta_{i_*G}\).  Hence
\(i_*(\varepsilon_G)\) is an equivalence.  By
\Ref{lem:closed-direct-image-conservative}, \(\varepsilon_G\) is an
equivalence.  Thus \(i_*\) is fully faithful.

Orthogonality shows that the essential image of \(i_*\) is contained in
\(\mathbf H_Z(S)\).  Conversely, if \(F\in\mathbf H_Z(S)\), then
\[
    j^*F
    \simeq
    *_U.
\]
The localization square therefore gives
\[
    F
    \xrightarrow{\simeq}
    i_*i^*F,
\]
so \(F\) belongs to the essential image.  This proves the unpointed
statement.  The pointed proof is identical, with terminal objects
replaced by zero objects.
\end{proof}

\begin{theorem}[Pointed localization fibre sequence]
\label{thm:pointed-localization-fiber-sequence}
For every
\[
    F
    \in
    \mathbf H_\bullet(S),
\]
there is a natural fibre sequence
\[
    i_*i^!F
    \longrightarrow
    F
    \longrightarrow
    j_*j^*F.
\]
\end{theorem}

\begin{proof}
Let
\[
    K
    :=
    \operatorname{fib}
    \bigl(
        F
        \longrightarrow
        j_*j^*F
    \bigr).
\]
Since
\[
    j_\sharp
    \dashv
    j^*,
\]
the functor \(j^*\) is a right adjoint and therefore preserves limits,
in particular fibres.  Applying \(j^*\), and using
\[
    j^*j_*
    \simeq
    \operatorname{id},
\]
gives
\[
\begin{aligned}
    j^*K
    &\simeq
    \operatorname{fib}
    \bigl(
        j^*F
        \longrightarrow
        j^*j_*j^*F
    \bigr)
    \\
    &\simeq
    \operatorname{fib}
    \bigl(
        j^*F
        \xrightarrow{\simeq}
        j^*F
    \bigr)
    \\
    &\simeq
    0.
\end{aligned}
\]
Thus \(K\) is supported on \(Z\).  By
\Ref{thm:closed-support-kashiwara},
there is a unique object
\[
    G
    \in
    \mathbf H_\bullet(Z)
\]
such that
\[
    K
    \simeq
    i_*G.
\]

For every
\[
    A
    \in
    \mathbf H_\bullet(Z),
\]
open--closed orthogonality gives
\[
\begin{aligned}
    \operatorname{Map}_S
    \bigl(
        i_*A,
        j_*j^*F
    \bigr)
    &\simeq
    \operatorname{Map}_U
    \bigl(
        j^*i_*A,
        j^*F
    \bigr)
    \\
    &\simeq
    *.
\end{aligned}
\]
Mapping \(i_*A\) into the defining fibre sequence for \(K\) therefore
gives
\[
\begin{aligned}
    \operatorname{Map}_Z(A,G)
    &\simeq
    \operatorname{Map}_S(i_*A,K)
    \\
    &\simeq
    \operatorname{Map}_S(i_*A,F)
    \\
    &\simeq
    \operatorname{Map}_Z(A,i^!F).
\end{aligned}
\]
Yoneda gives
\[
    G
    \simeq
    i^!F.
\]
Hence
\[
    K
    \simeq
    i_*i^!F,
\]
which is the asserted fibre sequence.
\end{proof}

\begin{corollary}[Concrete formula for closed exceptional pullback]
\label{cor:closed-exceptional-pullback-formula}
For
\[
    F
    \in
    \mathbf H_\bullet(S),
\]
there is a canonical equivalence
\[
    i^!F
    \simeq
    i^*
    \operatorname{fib}
    \bigl(
        F
        \longrightarrow
        j_*j^*F
    \bigr).
\]
Equivalently,
\[
    i_*i^!F
    \simeq
    \operatorname{fib}
    \bigl(
        F
        \longrightarrow
        j_*j^*F
    \bigr).
\]
\end{corollary}

\begin{proof}
This is the identification of the fibre in
\Ref{thm:pointed-localization-fiber-sequence}, together with the full
faithfulness of \(i_*\).
\end{proof}

\subsection{The recollement package}

\begin{theorem}[Closed--open recollement]
\label{thm:closed-open-recollement}
There is a recollement diagram
\[
    \mathbf H_\bullet(Z)
    \;
    \substack{
        \xrightarrow{\ i_*\ }
        \\
        \xleftarrow[\ i^*\ ]{}
        \\
        \xleftarrow[\ i^!\ ]{}
    }
    \;
    \mathbf H_\bullet(S)
    \;
    \substack{
        \xrightarrow{\ j^*\ }
        \\
        \xleftarrow[\ j_\sharp\ ]{}
        \\
        \xleftarrow[\ j_*\ ]{}
    }
    \;
    \mathbf H_\bullet(U).
\]
More precisely:
\begin{enumerate}[label=(\roman*)]
    \item \(i_*,j_\sharp,j_*\) are fully faithful;

    \item
    \[
        j^*i_*
        \simeq
        0,
        \qquad
        i^*j_\sharp
        \simeq
        0;
    \]

    \item the pair
    \[
        (i^*,j^*)
    \]
    is jointly conservative;

    \item every object is reconstructed by the cofibre and fibre
    sequences of
    \Ref{thm:pointed-localization-cofiber-sequence}
    and
    \Ref{thm:pointed-localization-fiber-sequence}.
\end{enumerate}
\end{theorem}

\begin{proof}
Full faithfulness of \(i_*\) is
\Ref{thm:closed-support-kashiwara}; full faithfulness of \(j_\sharp\)
and \(j_*\) is \Ref{prop:open-immersion-full-faithfulness}.
Orthogonality is \Ref{prop:open-closed-orthogonality}.

To prove joint conservativity on morphisms, let
\[
    \varphi
    \colon
    F
    \longrightarrow
    G
\]
be a morphism such that \(i^*\varphi\) and \(j^*\varphi\) are
equivalences.  Naturality gives a morphism from the pointed
localization pushout square of \(F\) to that of \(G\).  The induced maps
on the upper-left and lower-right vertices are
\[
    j_\sharp j^*\varphi
    \quad\text{and}\quad
    i_*i^*\varphi,
\]
and are equivalences.  The lower-left vertex is the zero object in both
squares.  Thus three corresponding vertices of the two cocartesian
squares are equivalent.  Universality of pushouts implies that the
remaining map \(\varphi\) is an equivalence.  Thus \((i^*,j^*)\) is
jointly conservative.  Reconstruction is exactly the pair of
localization sequences.
\end{proof}

\subsection{Closed base change}

\begin{theorem}[Closed base change]
\label{thm:unstable-closed-base-change}
Let
\[
\begin{tikzcd}[column sep=large, row sep=large]
    Y
        \arrow[r, "k"]
        \arrow[d, "g"']
    &
    X
        \arrow[d, "f"]
    \\
    Z
        \arrow[r, "i"']
    &
    S
\end{tikzcd}
\]
be a Cartesian square of quasi-compact quasi-separated spectral schemes,
where \(i\) is a closed immersion with quasi-compact open complement.
Then \(k\) is a closed immersion with quasi-compact open complement.

The canonical exchange transformation
\[
    \operatorname{Ex}^{*}_{*}
    \colon
    f^*i_*
    \longrightarrow
    k_*g^*
\]
is an equivalence on unpointed and pointed motivic spaces.
\end{theorem}

\begin{proof}
The exchange transformation is defined by the composite
\[
\begin{aligned}
    f^*i_*
    &\xrightarrow{\eta_k}
    k_*k^*f^*i_*
    \\
    &\xrightarrow{\simeq}
    k_*g^*i^*i_*
    \\
    &\xrightarrow{
        k_*g^*(\varepsilon_i)
    }
    k_*g^*,
\end{aligned}
\]
where
\[
    \eta_k
    \colon
    \operatorname{id}
    \longrightarrow
    k_*k^*
\]
is the unit of \(k^*\dashv k_*\), and
\[
    \varepsilon_i
    \colon
    i^*i_*
    \longrightarrow
    \operatorname{id}
\]
is the counit of \(i^*\dashv i_*\).

Let
\[
    j
    \colon
    U
    \hookrightarrow
    S
\]
be the open complement of \(i\), and let
\[
    j'
    \colon
    V
    :=
    X\times_SU
    \hookrightarrow
    X
\]
be its base change.  Then \(j'\) is the open complement of \(k\).

For \(F\in\mathbf H(Z)\), both
\[
    f^*i_*F
    \quad\text{and}\quad
    k_*g^*F
\]
are supported on \(Y\).  Indeed,
\[
\begin{aligned}
    (j')^*f^*i_*F
    &\simeq
    (f|_V)^*j^*i_*F
    \\
    &\simeq
    *_V,
\end{aligned}
\]
while
\[
    (j')^*k_*g^*F
    \simeq
    *_V
\]
by open--closed orthogonality.

Apply \(k^*\).  Using the Cartesian identity
\[
    k^*f^*
    \simeq
    g^*i^*
\]
and full faithfulness of \(i_*\) and \(k_*\), both sides identify with
\(g^*F\):
\[
\begin{aligned}
    k^*f^*i_*F
    &\simeq
    g^*i^*i_*F
    \\
    &\simeq
    g^*F,
    \\
    k^*k_*g^*F
    &\simeq
    g^*F.
\end{aligned}
\]
Under these identifications, the pullback of
\(\operatorname{Ex}^{*}_{*}\) is the identity of \(g^*F\), by the
triangle identities for the two adjunctions.

The support theorem identifies
\[
    k_*
    \colon
    \mathbf H(Y)
    \xrightarrow{\simeq}
    \mathbf H_Y(X).
\]
Therefore a morphism between objects supported on \(Y\) is an
equivalence if and only if its pullback by \(k^*\) is an equivalence.
It follows that
\[
    f^*i_*F
    \xrightarrow{\simeq}
    k_*g^*F.
\]

The pointed proof is identical.
\end{proof}

\begin{corollary}[Exceptional closed base change]
\label{cor:exceptional-closed-base-change}
In the Cartesian square of
\Ref{thm:unstable-closed-base-change},
there is a canonical equivalence
\[
    i^!f_*
    \simeq
    g_*k^!
\]
on pointed motivic spaces.
\end{corollary}

\begin{proof}
Take right adjoints of the closed base-change equivalence
\[
    f^*i_*
    \simeq
    k_*g^*.
\]
\end{proof}

\subsection{Closed projection formula}

\begin{theorem}[Closed projection formula]
\label{thm:unstable-closed-projection-formula}
For
\[
    F
    \in
    \mathbf H_\bullet(Z),
    \qquad
    G
    \in
    \mathbf H_\bullet(S),
\]
the canonical projection transformation is an equivalence:
\[
    i_*F
    \wedge_S
    G
    \xrightarrow{\simeq}
    i_*
    \bigl(
        F
        \wedge_Z
        i^*G
    \bigr).
\]
\end{theorem}

\begin{proof}
Both sides are supported on \(Z\).  The left side is supported because
\[
    j^*(i_*F\wedge_SG)
    \simeq
    j^*i_*F\wedge_Uj^*G
    \simeq
    0.
\]
The right side is supported because it lies in the image of \(i_*\).

Apply \(i^*\).  Strong symmetric monoidality of inverse image and
\(i^*i_*\simeq\id\) identify the left side with
\[
    F\wedge_Zi^*G,
\]
and the right side with the same object.  Under these identifications,
the projection transformation is the identity.  The support theorem
then implies that it is an equivalence.
\end{proof}

\begin{corollary}[Closed internal-Hom exchange]
\label{cor:closed-internal-hom-exchange}
For
\[
    G,H
    \in
    \mathbf H_\bullet(S),
\]
there is a canonical equivalence
\[
    i^!
    \underline{\operatorname{Map}}_{\mathbf H_\bullet(S)}(G,H)
    \simeq
    \underline{\operatorname{Map}}_{\mathbf H_\bullet(Z)}
    \bigl(
        i^*G,
        i^!H
    \bigr).
\]
\end{corollary}

\begin{proof}
Let
\[
    F
    \in
    \mathbf H_\bullet(Z).
\]
Using successively the adjunction \(i_*\dashv i^!\), the
smash--internal-Hom adjunction, the closed projection formula, and
again \(i_*\dashv i^!\), one obtains
\[
\begin{aligned}
    &
    \operatorname{Map}_Z
    \left(
        F,
        i^!
        \underline{\operatorname{Map}}_{\mathbf H_\bullet(S)}(G,H)
    \right)
    \\
    &\simeq
    \operatorname{Map}_S
    \left(
        i_*F,
        \underline{\operatorname{Map}}_{\mathbf H_\bullet(S)}(G,H)
    \right)
    \\
    &\simeq
    \operatorname{Map}_S
    \bigl(
        i_*F\wedge_SG,
        H
    \bigr)
    \\
    &\simeq
    \operatorname{Map}_S
    \left(
        i_*
        \bigl(
            F\wedge_Zi^*G
        \bigr),
        H
    \right)
    \\
    &\simeq
    \operatorname{Map}_Z
    \bigl(
        F\wedge_Zi^*G,
        i^!H
    \bigr)
    \\
    &\simeq
    \operatorname{Map}_Z
    \left(
        F,
        \underline{\operatorname{Map}}_{\mathbf H_\bullet(Z)}
        \bigl(
            i^*G,
            i^!H
        \bigr)
    \right).
\end{aligned}
\]
Yoneda proves the asserted equivalence.
\end{proof}

\subsection{Smooth--closed exchange}


\begin{theorem}[Pointed smooth--closed base change]
\label{thm:unstable-smooth-closed-base-change}
Let
\[
\begin{tikzcd}[column sep=large, row sep=large]
    Y
        \arrow[r, "k"]
        \arrow[d, "q"']
    &
    X
        \arrow[d, "p"]
    \\
    Z
        \arrow[r, "i"']
    &
    S
\end{tikzcd}
\]
be Cartesian, where \(p\) is smooth and of finite presentation and
\(i\) is a closed immersion with quasi-compact open complement.

The canonical exchange transformation
\[
    \operatorname{Ex}^{\sharp}_{*}
    \colon
    p_{\sharp,\bullet}k_*
    \longrightarrow
    i_*q_{\sharp,\bullet}
\]
is an equivalence on pointed motivic spaces.
\end{theorem}

\begin{proof}
The exchange transformation is defined by the composite
\[
\begin{aligned}
    p_{\sharp,\bullet}k_*
    &\xrightarrow{\eta_i}
    i_*i^*p_{\sharp,\bullet}k_*
    \\
    &\xrightarrow{\simeq}
    i_*q_{\sharp,\bullet}k^*k_*
    \\
    &\xrightarrow{
        i_*q_{\sharp,\bullet}(\varepsilon_k)
    }
    i_*q_{\sharp,\bullet}.
\end{aligned}
\]
Here the middle equivalence is pointed smooth base change
\[
    i^*p_{\sharp,\bullet}
    \simeq
    q_{\sharp,\bullet}k^*,
\]
and
\[
    \varepsilon_k
    \colon
    k^*k_*
    \longrightarrow
    \operatorname{id}
\]
is the counit of \(k^*\dashv k_*\).

Let
\[
    F
    \in
    \mathbf H_\bullet(Y).
\]
Both
\[
    p_{\sharp,\bullet}k_*F
    \quad\text{and}\quad
    i_*q_{\sharp,\bullet}F
\]
are supported on \(Z\).  The second assertion is immediate.

For the first, let
\[
    j
    \colon
    U
    \hookrightarrow
    S
\]
be the open complement of \(i\), and let
\[
    j'
    \colon
    X_U
    :=
    X\times_SU
    \hookrightarrow
    X
\]
be its base change.  Write
\[
    p_U
    \colon
    X_U
    \longrightarrow
    U
\]
for the base change of \(p\).  Pointed smooth base change gives
\[
\begin{aligned}
    j^*p_{\sharp,\bullet}k_*F
    &\simeq
    (p_U)_{\sharp,\bullet}(j')^*k_*F
    \\
    &\simeq
    (p_U)_{\sharp,\bullet}(0_{X_U})
    \\
    &\simeq
    0_U,
\end{aligned}
\]
where the second equivalence is open--closed orthogonality and the last
uses preservation of the zero object by the pointed left adjoint
\((p_U)_{\sharp,\bullet}\).  Thus
\[
    p_{\sharp,\bullet}k_*F
\]
is supported on \(Z\).

Apply \(i^*\).  Pointed smooth base change and full faithfulness of
\(k_*\) give
\[
\begin{aligned}
    i^*p_{\sharp,\bullet}k_*F
    &\simeq
    q_{\sharp,\bullet}k^*k_*F
    \\
    &\simeq
    q_{\sharp,\bullet}F.
\end{aligned}
\]
Likewise,
\[
    i^*i_*q_{\sharp,\bullet}F
    \simeq
    q_{\sharp,\bullet}F.
\]
Under these identifications, the pullback of
\(\operatorname{Ex}^{\sharp}_{*}\) is the identity of
\(q_{\sharp,\bullet}F\), by the triangle identity for
\(k^*\dashv k_*\).

Both the source and target of the exchange transformation are supported
on \(Z\), and the support theorem gives an equivalence
\[
    i_*
    \colon
    \mathbf H_\bullet(Z)
    \xrightarrow{\simeq}
    \mathbf H_{\bullet,Z}(S).
\]
A morphism between supported objects is therefore an equivalence if and
only if its pullback by \(i^*\) is an equivalence.  Consequently,
\[
    p_{\sharp,\bullet}k_*F
    \xrightarrow{\simeq}
    i_*q_{\sharp,\bullet}F.
\]
\end{proof}

\begin{corollary}[Exceptional smooth--closed exchange]
\label{cor:exceptional-smooth-closed-exchange}
In the Cartesian square of
\Ref{thm:unstable-smooth-closed-base-change},
there is a canonical equivalence
\[
    q^*i^!
    \simeq
    k^!p^*
\]
on pointed motivic spaces.
\end{corollary}

\begin{proof}
Take right adjoints of
\[
    p_\sharp k_*
    \simeq
    i_*q_\sharp.
\]
\end{proof}

\section{Derived nilpotent invariance and classical comparison}
\label{sec:derived-nilpotent-invariance-classical-comparison}

\subsection{Closed immersions with empty complement}

\begin{theorem}[Topological invariance for closed immersions]
\label{thm:topological-invariance-closed-immersion}
Let
\[
    i
    \colon
    S_0
    \hookrightarrow
    S
\]
be a closed immersion of quasi-compact quasi-separated spectral schemes
which induces a homeomorphism
\[
    |S_0|
    \xrightarrow{\simeq}
    |S|.
\]
Then inverse and direct image are mutually inverse equivalences
\[
    i^*
    \colon
    \mathbf H(S)
    \xrightarrow{\simeq}
    \mathbf H(S_0),
\]
and
\[
    i_*
    \colon
    \mathbf H(S_0)
    \xrightarrow{\simeq}
    \mathbf H(S).
\]
The same is true for pointed motivic spaces.
\end{theorem}

\begin{proof}
The open complement of \(S_0\) in \(S\) is empty.  The localization
square therefore gives
\[
    F
    \xrightarrow{\simeq}
    i_*i^*F
\]
for every \(F\in\mathbf H(S)\).  Full faithfulness of \(i_*\) gives
\[
    i^*i_*
    \simeq
    \operatorname{id}.
\]
Thus \(i^*\) and \(i_*\) are mutually inverse equivalences.  The
pointed assertion follows from pointed localization.
\end{proof}

\begin{corollary}[Invariance under equal reduced closed structure]
\label{cor:invariance-equal-reduced-structure}
Let
\[
    i
    \colon
    S_0
    \hookrightarrow
    S
\]
be a closed immersion inducing an equivalence on reduced classical
schemes:
\[
    (S_0)_{\mathrm{cl,red}}
    \xrightarrow{\simeq}
    S_{\mathrm{cl,red}}.
\]
Then
\[
    i^*
    \colon
    \mathbf H(S)
    \xrightarrow{\simeq}
    \mathbf H(S_0)
\]
and its pointed analogue are equivalences.
\end{corollary}

\begin{proof}
The two spectral schemes have the same underlying topological space.
Apply \Ref{thm:topological-invariance-closed-immersion}.
\end{proof}

\subsection{Derived nilpotent invariance}

\begin{theorem}[Derived nilpotent invariance]
\label{thm:derived-nilpotent-invariance}
For every quasi-compact quasi-separated spectral scheme \(S\), the
canonical closed immersion
\[
    \iota_S
    \colon
    S_{\mathrm{cl}}
    \longrightarrow
    S
\]
induces equivalences
\[
    \iota_S^*
    \colon
    \mathbf H(S)
    \xrightarrow{\simeq}
    \mathbf H(S_{\mathrm{cl}})
\]
and
\[
    \iota_S^*
    \colon
    \mathbf H_\bullet(S)
    \xrightarrow{\simeq}
    \mathbf H_\bullet(S_{\mathrm{cl}}).
\]
These equivalences are natural in \(S\).
\end{theorem}

\begin{proof}
The canonical map
\[
    S_{\mathrm{cl}}
    \longrightarrow
    S
\]
is a closed immersion and induces the identity on underlying
topological spaces.  Apply
\Ref{thm:topological-invariance-closed-immersion}.  Naturality follows
from functoriality of classical truncation and coherent base
functoriality of motivic inverse image.
\end{proof}

\begin{remark}[Meaning of derived nilpotent invariance]
\label{rem:meaning-derived-nilpotent-invariance}
The category
\[
    \mathbf H(S_{\mathrm{cl}})
\]
in \Ref{thm:derived-nilpotent-invariance} is still formed from
differentially smooth spectral schemes over the discrete spectral base
\(S_{\mathrm{cl}}\), using the brave new affine line.  Derived
nilpotent invariance alone does not identify it with the ordinary
Morel--Voevodsky category.  That comparison is the content of the next
subsection.
\end{remark}

\subsection{Comparison with ordinary motivic spaces}

\begin{definition}[Ordinary motivic spaces]
\label{def:ordinary-motivic-spaces}
Let \(S_0\) be an ordinary quasi-compact quasi-separated scheme.  Write
\[
    \mathbf H_{\mathrm{cl}}(S_0)
\]
for the ordinary unstable motivic \(\infty\)-category obtained from
smooth schemes of finite presentation over \(S_0\) by imposing ordinary
Nisnevich descent and ordinary affine-line invariance.  Its pointed
counterpart is denoted
\[
    \mathbf H_{\mathrm{cl},\bullet}(S_0).
\]
\end{definition}

\begin{construction}[Classical comparison adjunction]
\label{con:classical-motivic-comparison}
For a quasi-compact quasi-separated spectral scheme \(S\), classical
truncation of smooth test schemes defines a functor
\[
    w_S
    \colon
    \operatorname{Sm}^{\mathrm{fp}}_{/S}
    \longrightarrow
    \operatorname{Sm}^{\mathrm{fp},\mathrm{cl}}_{/S_{\mathrm{cl}}},
    \qquad
    X
    \longmapsto
    X_{\mathrm{cl}}.
\]
The local standard-form theorem implies that
\[
    X_{\mathrm{cl}}
    \longrightarrow
    S_{\mathrm{cl}}
\]
is ordinarily smooth and of finite presentation.

Write
\[
    \operatorname{PSh}_{\mathrm{cl}}(S_{\mathrm{cl}})
    :=
    \operatorname{Fun}
    \left(
        \bigl(
            \operatorname{Sm}^{\mathrm{fp},\mathrm{cl}}_{/S_{\mathrm{cl}}}
        \bigr)^{\mathrm{op}},
        \mathcal S
    \right).
\]
Left Kan extension along \(w_S^{\mathrm{op}}\) gives a
colimit-preserving functor
\[
    \widetilde\Theta_S
    :=
    \operatorname{Lan}_{w_S^{\mathrm{op}}}
    \colon
    \operatorname{PSh}(S)
    \longrightarrow
    \operatorname{PSh}_{\mathrm{cl}}(S_{\mathrm{cl}}).
\]
By the universal property of presheaf categories, it sends
representables to representables:
\[
    \widetilde\Theta_Sh_S(X)
    \simeq
    h^{\mathrm{cl}}_{S_{\mathrm{cl}}}
    \bigl(
        X_{\mathrm{cl}}
    \bigr).
\]

We verify that the composite
\[
    L^{\mathrm{cl}}_{\mathrm{mot},S_{\mathrm{cl}}}
    \widetilde\Theta_S
    \colon
    \operatorname{PSh}(S)
    \longrightarrow
    \mathbf H_{\mathrm{cl}}(S_{\mathrm{cl}})
\]
inverts every morphism in
\[
    W_{\mathrm{mot},S}
    =
    E_{\mathrm{Nis},S}
    \cup
    W_{\mathbb A^1,S}.
\]

Let
\[
\begin{tikzcd}[column sep=large, row sep=large]
    W
        \arrow[r]
        \arrow[d]
    &
    V
        \arrow[d]
    \\
    U
        \arrow[r]
    &
    X
\end{tikzcd}
\]
be a spectral Nisnevich distinguished square.  Classical truncation
preserves fibre products, open immersions, étale morphisms, and the
equivalence over the closed complement.  Hence
\[
\begin{tikzcd}[column sep=large, row sep=large]
    W_{\mathrm{cl}}
        \arrow[r]
        \arrow[d]
    &
    V_{\mathrm{cl}}
        \arrow[d]
    \\
    U_{\mathrm{cl}}
        \arrow[r]
    &
    X_{\mathrm{cl}}
\end{tikzcd}
\]
is an ordinary Nisnevich distinguished square.

It follows that \(\widetilde\Theta_S\) carries the corresponding
spectral excision morphism
\[
    h_S(U)
    \amalg_{h_S(W)}
    h_S(V)
    \longrightarrow
    h_S(X)
\]
to the ordinary excision morphism
\[
    h^{\mathrm{cl}}(U_{\mathrm{cl}})
    \amalg_{h^{\mathrm{cl}}(W_{\mathrm{cl}})}
    h^{\mathrm{cl}}(V_{\mathrm{cl}})
    \longrightarrow
    h^{\mathrm{cl}}(X_{\mathrm{cl}}).
\]
It also carries
\[
    \varnothing_{\operatorname{PSh}(S)}
    \longrightarrow
    h_S(\varnothing)
\]
to its ordinary analogue.  Therefore
\[
    L^{\mathrm{cl}}_{\mathrm{mot},S_{\mathrm{cl}}}
    \widetilde\Theta_S
\]
inverts every morphism in \(E_{\mathrm{Nis},S}\).

For every
\[
    X
    \in
    \operatorname{Sm}^{\mathrm{fp}}_{/S},
\]
classical truncation is compatible with the brave new affine line:
\[
\begin{aligned}
    \bigl(
        X\times_S\mathbb A^1_S
    \bigr)_{\mathrm{cl}}
    &\simeq
    X_{\mathrm{cl}}
    \times_{S_{\mathrm{cl}}}
    \bigl(
        \mathbb A^1_S
    \bigr)_{\mathrm{cl}}
    \\
    &\simeq
    X_{\mathrm{cl}}
    \times_{S_{\mathrm{cl}}}
    \mathbb A^1_{S_{\mathrm{cl}}}.
\end{aligned}
\]
Consequently, \(\widetilde\Theta_S\) carries the brave-new
affine-line projection
\[
    h_S
    \bigl(
        X\times_S\mathbb A^1_S
    \bigr)
    \longrightarrow
    h_S(X)
\]
to the ordinary affine-line projection
\[
    h^{\mathrm{cl}}_{S_{\mathrm{cl}}}
    \bigl(
        X_{\mathrm{cl}}
        \times_{S_{\mathrm{cl}}}
        \mathbb A^1_{S_{\mathrm{cl}}}
    \bigr)
    \longrightarrow
    h^{\mathrm{cl}}_{S_{\mathrm{cl}}}
    \bigl(
        X_{\mathrm{cl}}
    \bigr).
\]
Thus
\[
    L^{\mathrm{cl}}_{\mathrm{mot},S_{\mathrm{cl}}}
    \widetilde\Theta_S
\]
also inverts every morphism in \(W_{\mathbb A^1,S}\).

By the universal property of
\[
    L_{\mathrm{mot},S}
    \colon
    \operatorname{PSh}(S)
    \longrightarrow
    \mathbf H(S),
\]
there is therefore a unique colimit-preserving functor
\[
    \Theta_S
    \colon
    \mathbf H(S)
    \longrightarrow
    \mathbf H_{\mathrm{cl}}(S_{\mathrm{cl}})
\]
equipped with a canonical equivalence
\[
    \Theta_SL_{\mathrm{mot},S}
    \simeq
    L^{\mathrm{cl}}_{\mathrm{mot},S_{\mathrm{cl}}}
    \widetilde\Theta_S.
\]

Since \(\Theta_S\) is colimit-preserving between presentable
\(\infty\)-categories, it admits a right adjoint
\[
    \Psi_S
    \colon
    \mathbf H_{\mathrm{cl}}(S_{\mathrm{cl}})
    \longrightarrow
    \mathbf H(S).
\]

For every smooth spectral \(S\)-scheme \(X\) of finite presentation,
there is a canonical equivalence
\[
\begin{aligned}
    \Theta_S\mathrm M_S(X)
    &=
    \Theta_SL_{\mathrm{mot},S}h_S(X)
    \\
    &\simeq
    L^{\mathrm{cl}}_{\mathrm{mot},S_{\mathrm{cl}}}
    \widetilde\Theta_Sh_S(X)
    \\
    &\simeq
    L^{\mathrm{cl}}_{\mathrm{mot},S_{\mathrm{cl}}}
    h^{\mathrm{cl}}_{S_{\mathrm{cl}}}
    \bigl(
        X_{\mathrm{cl}}
    \bigr)
    \\
    &=
    \mathrm M^{\mathrm{cl}}_{S_{\mathrm{cl}}}
    \bigl(
        X_{\mathrm{cl}}
    \bigr).
\end{aligned}
\]

The fact that this adjunction is an equivalence, together with its
compatibility with base change, Cartesian products, pointing, and the
\((*,\sharp,\times)\)-formalism, is established in
\Ref{thm:classical-comparison-brave-new-motivic-spaces}.
\end{construction}

\begin{theorem}[Comparison with ordinary motivic homotopy theory]
\label{thm:classical-comparison-brave-new-motivic-spaces}
For every quasi-compact quasi-separated spectral scheme \(S\), the
classical comparison adjunction
\[
    \Theta_S
    \colon
    \mathbf H(S)
    \rightleftarrows
    \mathbf H_{\mathrm{cl}}(S_{\mathrm{cl}})
    :
    \Psi_S
\]
is an adjoint equivalence.

It induces a strong smash-monoidal equivalence
\[
    \Theta_{S,\bullet}
    \colon
    \mathbf H_\bullet(S)
    \xrightarrow{\simeq}
    \mathbf H_{\mathrm{cl},\bullet}(S_{\mathrm{cl}}).
\]

These equivalences are natural in \(S\).  Under them:
\begin{enumerate}[label=(\roman*)]
    \item inverse image corresponds to ordinary inverse image;

    \item direct image corresponds to ordinary direct image;

    \item smooth direct image corresponds to ordinary smooth direct
    image;

    \item Cartesian products and smash products correspond to their
    ordinary counterparts;

    \item the corresponding Cartesian and pointed internal mapping
    objects are preserved.
\end{enumerate}
\end{theorem}

\begin{proof}
The comparison theorem cited below is formulated for motivic
categories constructed from smooth spectral algebraic spaces rather
than only from smooth spectral schemes.  Since \(S\) is a spectral
scheme,
\cite[Corollary~2.4.5]{KhanBraveNewMotivicI}
shows that restriction along the inclusion
\[
    \operatorname{Sm}^{\mathrm{fp}}_{/S}
    \hookrightarrow
    \operatorname{Sm}^{\mathrm{fp},\mathrm{asp}}_{/S}
\]
induces an equivalence between the corresponding categories of
\(\mathbb A^1\)-invariant Nisnevich sheaves.

We require the analogous comparison on the ordinary side.  Let
\[
    \operatorname{Sm}^{\mathrm{fp},\mathrm{cl}}_{/S_{\mathrm{cl}}}
    \hookrightarrow
    \operatorname{Sm}^{\mathrm{fp},\mathrm{cl},\mathrm{asp}}_
    {/S_{\mathrm{cl}}}
\]
be the inclusion of smooth ordinary schemes into smooth ordinary
algebraic spaces of finite presentation over \(S_{\mathrm{cl}}\).
Every smooth ordinary algebraic space of finite presentation admits a
Nisnevich covering by smooth schemes of finite presentation.  Moreover,
fibre products of smooth schemes over such an algebraic space admit
Nisnevich coverings by smooth schemes.  Thus smooth schemes form a
Nisnevich basis for the smooth algebraic-space site.

The comparison lemma for a basis therefore gives an equivalence between
the corresponding categories of ordinary Nisnevich sheaves.
The inclusion is compatible with products by
\[
    \mathbb A^1_{S_{\mathrm{cl}}},
\]
so this equivalence restricts to an equivalence between the
\(\mathbb A^1\)-invariant Nisnevich sheaf categories.  Consequently, the
ordinary motivic category formed from smooth schemes agrees with the
one formed from smooth algebraic spaces.

Under these restriction equivalences, the comparison left adjoint of
\cite[Corollary~2.7.5]{CisinskiKhanA1Invariance}
agrees with the functor
\[
    \Theta_S
    \colon
    \mathbf H(S)
    \longrightarrow
    \mathbf H_{\mathrm{cl}}(S_{\mathrm{cl}})
\]
constructed above.  Indeed, both functors preserve colimits and carry
every smooth spectral motive to the ordinary motive of its classical
truncation:
\[
    \mathrm M_S(X)
    \longmapsto
    \mathrm M^{\mathrm{cl}}_{S_{\mathrm{cl}}}
    \bigl(
        X_{\mathrm{cl}}
    \bigr).
\]
The smooth scheme motives generate the scheme-based category under
sifted colimits by
\Ref{prop:generation-by-standard-smooth-motives}.
Thus the two comparison functors are canonically equivalent.

It now follows from
\cite[Corollary~2.7.5]{CisinskiKhanA1Invariance}
that the adjunction
\[
    \Theta_S
    \colon
    \mathbf H(S)
    \rightleftarrows
    \mathbf H_{\mathrm{cl}}(S_{\mathrm{cl}})
    :
    \Psi_S
\]
is an adjoint equivalence.

The comparison equivalences assemble pseudonaturally in the base.  Let
\[
    f
    \colon
    T
    \longrightarrow
    S
\]
be a morphism of quasi-compact quasi-separated spectral schemes.
Compatibility of classical truncation with geometric base change
gives a canonical equivalence
\[
    \Theta_Tf^*
    \simeq
    f_{\mathrm{cl}}^*\Theta_S.
\]
This proves compatibility with inverse image.

Taking right adjoints of this equivalence and using that
\[
    \Theta_S
    \qquad\text{and}\qquad
    \Theta_T
\]
are equivalences gives a canonical equivalence
\[
    \Theta_Sf_*
    \simeq
    (f_{\mathrm{cl}})_*\Theta_T.
\]
Thus direct image corresponds to ordinary direct image.

Now suppose that
\[
    p
    \colon
    X
    \longrightarrow
    S
\]
is smooth and of finite presentation.  The inverse-image comparison
\[
    \Theta_Xp^*
    \simeq
    p_{\mathrm{cl}}^*\Theta_S
\]
is an equivalence between right adjoints.  Passing to the corresponding
left adjoints gives
\[
    \Theta_Sp_\sharp
    \simeq
    (p_{\mathrm{cl}})_\sharp\Theta_X.
\]
Hence smooth direct image corresponds to ordinary smooth direct image.

These compatibilities agree with the functorial comparison described
in
\cite[Remark~2.7.9]{CisinskiKhanA1Invariance}.

We next verify the Cartesian symmetric monoidal compatibility.  On
smooth motives one has
\[
\begin{aligned}
    \Theta_S
    \bigl(
        \mathrm M_S(X)
        \times
        \mathrm M_S(Y)
    \bigr)
    &\simeq
    \Theta_S
    \mathrm M_S(X\times_SY)
    \\
    &\simeq
    \mathrm M^{\mathrm{cl}}_{S_{\mathrm{cl}}}
    \bigl(
        X_{\mathrm{cl}}
        \times_{S_{\mathrm{cl}}}
        Y_{\mathrm{cl}}
    \bigr)
    \\
    &\simeq
    \Theta_S\mathrm M_S(X)
    \times
    \Theta_S\mathrm M_S(Y).
\end{aligned}
\]
The bifunctors
\[
    (F,G)
    \longmapsto
    \Theta_S(F\times G)
\]
and
\[
    (F,G)
    \longmapsto
    \Theta_S(F)\times\Theta_S(G)
\]
preserve colimits separately in \(F\) and \(G\).  Since smooth motives
generate \(\mathbf H(S)\) under colimits, the displayed comparison
extends uniquely to a canonical equivalence
\[
    \Theta_S(F\times G)
    \simeq
    \Theta_S(F)\times\Theta_S(G)
\]
for arbitrary
\[
    F,G
    \in
    \mathbf H(S).
\]
The terminal motive is also preserved.  Therefore \(\Theta_S\) is
strong Cartesian symmetric monoidal.

Since \(\Theta_S\) preserves the terminal object, it induces an
equivalence of pointed categories
\[
    \Theta_{S,\bullet}
    \colon
    \mathbf H_\bullet(S)
    \xrightarrow{\simeq}
    \mathbf H_{\mathrm{cl},\bullet}(S_{\mathrm{cl}}).
\]
The smash product is constructed from Cartesian products and colimits
by
\[
    F\wedge G
    =
    (F\times G)
    \amalg_{F\amalg G}
    *,
\]
so the pointed comparison equivalence is strong symmetric monoidal for
the smash products.

Finally, a strong monoidal equivalence between closed symmetric
monoidal \(\infty\)-categories preserves internal mapping objects.
Thus there are canonical equivalences
\[
    \Theta_S
    \underline{\operatorname{Map}}_{\mathbf H(S)}
    (F,G)
    \simeq
    \underline{\operatorname{Map}}_{
        \mathbf H_{\mathrm{cl}}(S_{\mathrm{cl}})
    }
    \bigl(
        \Theta_SF,
        \Theta_SG
    \bigr)
\]
and
\[
    \Theta_{S,\bullet}
    \underline{\operatorname{Map}}_{\mathbf H_\bullet(S)}
    (F,G)
    \simeq
    \underline{\operatorname{Map}}_{
        \mathbf H_{\mathrm{cl},\bullet}(S_{\mathrm{cl}})
    }
    \bigl(
        \Theta_{S,\bullet}F,
        \Theta_{S,\bullet}G
    \bigr).
\]
This proves all the asserted compatibilities.
\end{proof}

\section{Homotopy purity}
\label{sec:brave-new-homotopy-purity}

\subsection{Smooth closed pairs and Thom spaces}

\begin{definition}[Smooth closed pair]
\label{def:smooth-closed-pair}
A \textbf{smooth closed pair over \(S\)} is a closed immersion
\[
    i
    \colon
    Z
    \hookrightarrow
    X
\]
such that both
\[
    X
    \longrightarrow
    S
\]
and
\[
    Z
    \longrightarrow
    S
\]
are smooth and of finite presentation.
\end{definition}

By \Ref{thm:closed-between-smooth-regular}, a smooth closed pair is a
regular closed immersion.  Its conormal module
\[
    N^*_{Z/X}
    :=
    L_{Z/X}[-1]
\]
is finite locally free.

\begin{notation}[Conormal module, normal module, and normal bundle]
\label{not:normal-data-smooth-closed-pair}
For a smooth closed pair
\[
    i
    \colon
    Z
    \hookrightarrow
    X,
\]
write
\[
    N^*_{Z/X}
    :=
    L_{Z/X}[-1]
\]
for the conormal module and
\[
    N_{Z/X}
    :=
    \bigl(
        N^*_{Z/X}
    \bigr)^\vee
\]
for the normal module.

The normal bundle is the relative linear space
\[
    \mathbb N_{Z/X}
    :=
    \mathbb V_Z
    \bigl(
        N^*_{Z/X}
    \bigr).
\]
Since
\[
    N_{Z/X}
    =
    \bigl(
        N^*_{Z/X}
    \bigr)^\vee,
\]
finite-projective biduality gives a canonical equivalence
\[
    N_{Z/X}^\vee
    \simeq
    N^*_{Z/X}.
\]
Consequently,
\[
    \mathbb N_{Z/X}
    \simeq
    \mathbb V_Z
    \bigl(
        N_{Z/X}^\vee
    \bigr),
\]
so \(\mathbb N_{Z/X}\) is the geometric vector bundle associated with
the normal module \(N_{Z/X}\) under the convention
\[
    \mathbb V_Z(E)
    =
    \operatorname{Spec}_Z
    \operatorname{Sym}_{\mathcal O_Z}(E).
\]
Its zero section is denoted
\[
    s
    \colon
    Z
    \hookrightarrow
    \mathbb N_{Z/X}.
\]
\end{notation}

\begin{definition}[Unstable Thom space]
\label{def:unstable-thom-space}
Let \(E\) be a finite locally free quasi-coherent module on a smooth
spectral \(S\)-scheme \(Z\).  Define
\[
    \operatorname{Th}_S(E)
    :=
    \mathbb V_Z(E)/
    \bigl(
        \mathbb V_Z(E)\setminus Z
    \bigr)
    \in
    \mathbf H_\bullet(S),
\]
where \(Z\) is embedded by the zero section and both schemes are viewed
over \(S\).

For a smooth closed pair
\[
    Z
    \hookrightarrow
    X,
\]
define the Thom space of its normal bundle by
\[
\begin{aligned}
    \operatorname{Th}_S
    \bigl(
        \mathbb N_{Z/X}
    \bigr)
    &:=
    \mathbb N_{Z/X}/
    \bigl(
        \mathbb N_{Z/X}\setminus Z
    \bigr)
    \\
    &=
    \operatorname{Th}_S
    \bigl(
        N^*_{Z/X}
    \bigr).
\end{aligned}
\]
\end{definition}

\subsection{Nisnevich excision for pairs}


\begin{proposition}[Nisnevich excision for closed pairs]
\label{prop:nisnevich-excision-closed-pairs}
Let \(S\) be a quasi-compact quasi-separated spectral scheme.  Let
\[
    X
    \longrightarrow
    S
\]
be smooth and of finite presentation, and let
\[
    Z
    \hookrightarrow
    X
\]
be a closed immersion with quasi-compact open complement
\[
    U
    :=
    X\setminus Z.
\]

Let
\[
    q
    \colon
    X'
    \longrightarrow
    X
\]
be étale and of finite presentation.  Put
\[
    Z'
    :=
    X'\times_XZ,
    \qquad
    U'
    :=
    X'\times_XU,
\]
and suppose that the projection
\[
    Z'
    \longrightarrow
    Z
\]
is an equivalence.  Then the canonical morphism
\[
    X'/U'
    \longrightarrow
    X/U
\]
is an equivalence in
\[
    \mathbf H_\bullet(S).
\]
\end{proposition}

\begin{proof}
Because \(q\) is étale and of finite presentation and
\(X\to S\) is smooth and of finite presentation, the composite
\[
    X'
    \longrightarrow
    S
\]
is smooth and of finite presentation.  The open immersions
\[
    U
    \hookrightarrow
    X
    \qquad\text{and}\qquad
    U'
    \hookrightarrow
    X'
\]
are quasi-compact and therefore of finite presentation.  Consequently,
\(U\) and \(U'\) also belong to
\[
    \operatorname{Sm}^{\mathrm{fp}}_{/S}.
\]

The Cartesian square
\[
\begin{tikzcd}[column sep=large, row sep=large]
    U'
        \arrow[r]
        \arrow[d]
    &
    X'
        \arrow[d, "q"]
    \\
    U
        \arrow[r]
    &
    X
\end{tikzcd}
\]
is a spectral Nisnevich distinguished square: the lower horizontal
morphism is an open immersion, the right vertical morphism is étale,
and the induced morphism over the closed complement is the assumed
equivalence
\[
    Z'
    \xrightarrow{\simeq}
    Z.
\]

Applying the motivic Yoneda functor gives a cocartesian square
\[
\begin{tikzcd}[column sep=large, row sep=large]
    \mathrm M_S(U')
        \arrow[r]
        \arrow[d]
    &
    \mathrm M_S(X')
        \arrow[d]
    \\
    \mathrm M_S(U)
        \arrow[r]
    &
    \mathrm M_S(X)
\end{tikzcd}
\]
in \(\mathbf H(S)\).

Push out this square along the terminal morphism
\[
    \mathrm M_S(U)
    \longrightarrow
    *_S.
\]
By associativity of pushouts, the resulting comparison identifies the
two pointed quotients:
\[
    \mathrm M_S(X')
    \amalg_{\mathrm M_S(U')}
    *_S
    \xrightarrow{\simeq}
    \mathrm M_S(X)
    \amalg_{\mathrm M_S(U)}
    *_S.
\]
These are precisely
\[
    X'/U'
    \qquad\text{and}\qquad
    X/U.
\]
\end{proof}

\subsection{Classical comparison of normal data}

\begin{proposition}[Classical conormal comparison]
\label{prop:classical-conormal-comparison}
Let
\[
    i
    \colon
    Z
    \hookrightarrow
    X
\]
be a smooth closed pair over \(S\), and let \(\mathcal I\) be the ideal
of the ordinary closed immersion
\[
    i_{\mathrm{cl}}
    \colon
    Z_{\mathrm{cl}}
    \hookrightarrow
    X_{\mathrm{cl}}.
\]
There is a canonical equivalence of finite locally free
\(\mathcal O_{Z_{\mathrm{cl}}}\)-modules
\[
    N^*_{Z/X}
    \otimes_{\mathcal O_Z}
    \mathcal O_{Z_{\mathrm{cl}}}
    \simeq
    \mathcal I/\mathcal I^2.
\]
Consequently,
\[
    (\mathbb N_{Z/X})_{\mathrm{cl}}
    \simeq
    \mathbb V_{Z_{\mathrm{cl}}}
    (\mathcal I/\mathcal I^2),
\]
the ordinary normal bundle of the classical closed pair.
\end{proposition}

\begin{proof}
The transitivity triangle for
\[
    Z
    \xrightarrow{i}
    X
    \longrightarrow
    S
\]
is
\[
    i^*L_{X/S}
    \longrightarrow
    L_{Z/S}
    \longrightarrow
    L_{Z/X}.
\]
Since
\[
    L_{Z/X}
    \simeq
    N^*_{Z/X}[1],
\]
rotation gives a fibre sequence
\[
    N^*_{Z/X}
    \longrightarrow
    i^*L_{X/S}
    \longrightarrow
    L_{Z/S}.
\]

Tensor this fibre sequence with
\[
    \mathcal O_{Z_{\mathrm{cl}}}.
\]
Because
\[
    X\longrightarrow S,
    \qquad
    Z\longrightarrow S,
\]
are smooth, the modules
\[
    L_{X/S},
    \qquad
    L_{Z/S}
\]
are finite locally free.  The conormal module
\[
    N^*_{Z/X}
\]
is also finite locally free by
\Ref{thm:closed-between-smooth-regular}.  Their base changes to the
discrete spectral scheme \(Z_{\mathrm{cl}}\) are therefore discrete
finite locally free modules.

The degree-zero comparison for cotangent complexes gives canonical
isomorphisms
\[
    \pi_0
    \left(
        L_{X/S}
        \otimes_{\mathcal O_X}
        \mathcal O_{X_{\mathrm{cl}}}
    \right)
    \simeq
    \Omega^1_{X_{\mathrm{cl}}/S_{\mathrm{cl}}}
\]
and
\[
    \pi_0
    \left(
        L_{Z/S}
        \otimes_{\mathcal O_Z}
        \mathcal O_{Z_{\mathrm{cl}}}
    \right)
    \simeq
    \Omega^1_{Z_{\mathrm{cl}}/S_{\mathrm{cl}}}.
\]
Since the two base-changed cotangent complexes are discrete, these
degree-zero identifications determine equivalences
\[
    L_{X/S}
    \otimes_{\mathcal O_X}
    \mathcal O_{X_{\mathrm{cl}}}
    \simeq
    \Omega^1_{X_{\mathrm{cl}}/S_{\mathrm{cl}}}
\]
and
\[
    L_{Z/S}
    \otimes_{\mathcal O_Z}
    \mathcal O_{Z_{\mathrm{cl}}}
    \simeq
    \Omega^1_{Z_{\mathrm{cl}}/S_{\mathrm{cl}}},
\]
where the ordinary modules on the right are regarded as discrete
quasi-coherent spectral modules.

It follows that the base-changed fibre sequence is a fibre sequence of
discrete modules
\[
    N^*_{Z/X}
    \otimes_{\mathcal O_Z}
    \mathcal O_{Z_{\mathrm{cl}}}
    \longrightarrow
    i_{\mathrm{cl}}^*
    \Omega^1_{X_{\mathrm{cl}}/S_{\mathrm{cl}}}
    \longrightarrow
    \Omega^1_{Z_{\mathrm{cl}}/S_{\mathrm{cl}}}.
\]
Passing to degree zero gives a short exact sequence
\[
    0
    \longrightarrow
    \pi_0
    \left(
        N^*_{Z/X}
        \otimes_{\mathcal O_Z}
        \mathcal O_{Z_{\mathrm{cl}}}
    \right)
    \longrightarrow
    i_{\mathrm{cl}}^*
    \Omega^1_{X_{\mathrm{cl}}/S_{\mathrm{cl}}}
    \longrightarrow
    \Omega^1_{Z_{\mathrm{cl}}/S_{\mathrm{cl}}}
    \longrightarrow
    0.
\]

The ordinary conormal sequence for the regular closed immersion
\[
    Z_{\mathrm{cl}}
    \hookrightarrow
    X_{\mathrm{cl}}
\]
is
\[
    0
    \longrightarrow
    \mathcal I/\mathcal I^2
    \longrightarrow
    i_{\mathrm{cl}}^*
    \Omega^1_{X_{\mathrm{cl}}/S_{\mathrm{cl}}}
    \longrightarrow
    \Omega^1_{Z_{\mathrm{cl}}/S_{\mathrm{cl}}}
    \longrightarrow
    0.
\]
The maps to the last two terms are the same under the canonical
cotangent comparisons.  Therefore their fibres are canonically
equivalent:
\[
    N^*_{Z/X}
    \otimes_{\mathcal O_Z}
    \mathcal O_{Z_{\mathrm{cl}}}
    \simeq
    \mathcal I/\mathcal I^2.
\]

For a connective commutative algebra \(A\) and a finite projective
\(A\)-module \(P\), the universal property of the free commutative
algebra gives
\[
    \pi_0
    \operatorname{Sym}_A(P)
    \simeq
    \operatorname{Sym}_{\pi_0(A)}
    \pi_0(P).
\]
Applying this affine-locally to \(N^*_{Z/X}\), using the conormal
comparison just proved, and then gluing gives
\[
\begin{aligned}
    (\mathbb N_{Z/X})_{\mathrm{cl}}
    &=
    \left(
        \mathbb V_Z(N^*_{Z/X})
    \right)_{\mathrm{cl}}
    \\
    &\simeq
    \mathbb V_{Z_{\mathrm{cl}}}
    \left(
        N^*_{Z/X}
        \otimes_{\mathcal O_Z}
        \mathcal O_{Z_{\mathrm{cl}}}
    \right)
    \\
    &\simeq
    \mathbb V_{Z_{\mathrm{cl}}}
    (\mathcal I/\mathcal I^2).
\end{aligned}
\]
Under the convention
\[
    \mathbb V(E)
    =
    \operatorname{Spec}\operatorname{Sym}(E),
\]
this is the geometric vector bundle associated with the ordinary normal
module
\[
    (\mathcal I/\mathcal I^2)^\vee.
\]
\end{proof}

\begin{proposition}[Classical comparison of quotients and Thom spaces]
\label{prop:classical-comparison-quotient-thom}
For a smooth closed pair
\[
    Z
    \hookrightarrow
    X
\]
over \(S\), the pointed comparison equivalence of
\Ref{thm:classical-comparison-brave-new-motivic-spaces}
induces canonical equivalences
\[
    \Theta_{S,\bullet}
    \bigl(
        X/(X\setminus Z)
    \bigr)
    \simeq
    X_{\mathrm{cl}}/
    \bigl(
        X_{\mathrm{cl}}\setminus Z_{\mathrm{cl}}
    \bigr)
\]
and
\[
    \Theta_{S,\bullet}
    \operatorname{Th}_S
    \bigl(
        \mathbb N_{Z/X}
    \bigr)
    \simeq
    \operatorname{Th}^{\mathrm{cl}}_{S_{\mathrm{cl}}}
    \bigl(
        N_{Z_{\mathrm{cl}}/X_{\mathrm{cl}}}
    \bigr).
\]
\end{proposition}

\begin{proof}
The comparison functor is colimit-preserving and sends the motive of a
smooth spectral scheme to the ordinary motive of its classical
truncation.  Classical truncation preserves the open complement of a
closed immersion.  It therefore sends the pushout defining
\[
    X/(X\setminus Z)
\]
to the ordinary pointed quotient
\[
    X_{\mathrm{cl}}/
    \bigl(
        X_{\mathrm{cl}}\setminus Z_{\mathrm{cl}}
    \bigr).
\]

By
\Ref{prop:classical-conormal-comparison},
there is a canonical equivalence
\[
    (\mathbb N_{Z/X})_{\mathrm{cl}}
    \simeq
    N_{Z_{\mathrm{cl}}/X_{\mathrm{cl}}},
\]
where the right side denotes the ordinary normal bundle.  Classical
truncation also preserves the complement of the zero section.  Applying
the comparison functor to the pushout defining
\[
    \operatorname{Th}_S
    \bigl(
        \mathbb N_{Z/X}
    \bigr)
\]
therefore gives the ordinary Thom quotient.
\end{proof}

\subsection{Homotopy purity by classical comparison}


\begin{definition}[Étale morphism of smooth closed pairs]
\label{def:etale-morphism-smooth-closed-pairs}
Let \(S\) be a spectral scheme, and let
\[
    i
    \colon
    Z
    \hookrightarrow
    X
\]
and
\[
    i'
    \colon
    Z'
    \hookrightarrow
    X'
\]
be smooth closed pairs over \(S\).

An \textbf{étale morphism of smooth closed pairs}
\[
    (X',Z')
    \longrightarrow
    (X,Z)
\]
is a Cartesian square of spectral \(S\)-schemes
\[
\begin{tikzcd}[column sep=large, row sep=large]
    Z'
        \arrow[r]
        \arrow[d, hook, "i'"']
    &
    Z
        \arrow[d, hook, "i"]
    \\
    X'
        \arrow[r, "q"']
    &
    X
\end{tikzcd}
\]
in which
\[
    q
    \colon
    X'
    \longrightarrow
    X
\]
is étale.

When \(S\), \(X\), \(X'\), \(Z\), and \(Z'\) are ordinary schemes,
this recovers the ordinary notion through the fully faithful discrete
enhancement of ordinary schemes as spectral schemes.
\end{definition}


\begin{lemma}[Arbitrary base change for deformation spaces]
\label{lem:deformation-space-arbitrary-base-change}
Let \(S_0\) be an ordinary scheme, and let
\[
    Z_0
    \hookrightarrow
    X_0
\]
be a closed immersion such that both
\[
    X_0
    \longrightarrow
    S_0
    \qquad\text{and}\qquad
    Z_0
    \longrightarrow
    S_0
\]
are smooth and of finite presentation.  Let
\[
    f
    \colon
    T_0
    \longrightarrow
    S_0
\]
be an arbitrary morphism, and put
\[
    X_T
    :=
    X_0\times_{S_0}T_0,
    \qquad
    Z_T
    :=
    Z_0\times_{S_0}T_0.
\]

Then there is a canonical equivalence
\[
    D_{Z_0}X_0
    \times_{S_0}
    T_0
    \xrightarrow{\simeq}
    D_{Z_T}X_T
\]
over
\[
    \mathbb A^1_{T_0}.
\]
It is compatible with the closed immersions
\[
    Z_0\times\mathbb A^1
    \hookrightarrow
    D_{Z_0}X_0
\]
and
\[
    Z_T\times\mathbb A^1
    \hookrightarrow
    D_{Z_T}X_T,
\]
and with the canonical identifications of the fibres over \(1\) and
\(0\) with \(X_0,X_T\) and their respective normal bundles.
\end{lemma}

\begin{proof}
The assertion is local on \(S_0\) and \(X_0\).  Write
\[
    S_0
    =
    \Spec(R),
    \qquad
    X_0
    =
    \Spec(A),
    \qquad
    Z_0
    =
    \Spec(A/I).
\]
Since \(X_0\to S_0\) and \(Z_0\to S_0\) are smooth, both
\[
    A
    \qquad\text{and}\qquad
    A/I
\]
are flat over \(R\).

The immersion \(Z_0\hookrightarrow X_0\) is regular.  Consequently,
\[
    \operatorname{gr}_I(A)
    \simeq
    \operatorname{Sym}_{A/I}(I/I^2),
\]
and \(I/I^2\) is finite locally free over \(A/I\).  In particular, every
graded piece
\[
    I^n/I^{n+1}
    \simeq
    \operatorname{Sym}_{A/I}^n(I/I^2)
\]
is flat over \(R\).

The exact sequence
\[
    0
    \longrightarrow
    I
    \longrightarrow
    A
    \longrightarrow
    A/I
    \longrightarrow
    0
\]
shows that \(I\) is \(R\)-flat.  Inductively, the exact sequences
\[
    0
    \longrightarrow
    I^{n+1}
    \longrightarrow
    I^n
    \longrightarrow
    I^n/I^{n+1}
    \longrightarrow
    0
\]
show that every \(I^n\) is \(R\)-flat.

Let
\[
    R
    \longrightarrow
    R'
\]
be the ring map corresponding to the restriction of \(f\) to an affine
open of \(T_0\), and put
\[
    A'
    :=
    A\otimes_RR'.
\]
Since \(A/I\) is \(R\)-flat, tensoring the defining sequence for \(I\)
gives an exact sequence
\[
    0
    \longrightarrow
    I\otimes_RR'
    \longrightarrow
    A'
    \longrightarrow
    (A/I)\otimes_RR'
    \longrightarrow
    0.
\]
Thus the ideal of \(Z_T\) in \(X_T\) is
\[
    I'
    :=
    I\otimes_RR'
    \subseteq
    A'.
\]

The base-changed pair is again a closed immersion between schemes
smooth over \(T_0\), and hence is regular.  Conormal base change gives
\[
    I'/I'^2
    \simeq
    (I/I^2)\otimes_RR'.
\]
Therefore
\[
\begin{aligned}
    I'^n/I'^{n+1}
    &\simeq
    \operatorname{Sym}_{A'/I'}^n(I'/I'^2)
    \\
    &\simeq
    \operatorname{Sym}_{A/I}^n(I/I^2)
    \otimes_RR'
    \\
    &\simeq
    (I^n/I^{n+1})\otimes_RR'.
\end{aligned}
\]

We prove inductively that the natural morphisms
\[
    I^n\otimes_RR'
    \longrightarrow
    I'^n
\]
are isomorphisms.  For \(n=0\), this is the identity
\[
    A\otimes_RR'
    \simeq
    A'.
\]
Suppose it is known for \(n\).  Since \(I^n/I^{n+1}\) is \(R\)-flat,
tensoring the exact sequence
\[
    0
    \longrightarrow
    I^{n+1}
    \longrightarrow
    I^n
    \longrightarrow
    I^n/I^{n+1}
    \longrightarrow
    0
\]
with \(R'\) remains exact.  Comparing it with
\[
    0
    \longrightarrow
    I'^{n+1}
    \longrightarrow
    I'^n
    \longrightarrow
    I'^n/I'^{n+1}
    \longrightarrow
    0
\]
and using the induction hypothesis and the preceding identification of
the quotients gives
\[
    I^{n+1}\otimes_RR'
    \xrightarrow{\simeq}
    I'^{n+1}.
\]
The induction is complete.

The extended Rees algebra of \(I\) is
\[
    \mathcal R_I^{\mathrm{ext}}(A)
    :=
    A[t,It^{-1}]
    =
    \bigoplus_{n\in\mathbb Z}
    I^n t^{-n}
    \subseteq
    A[t,t^{-1}],
\]
where
\[
    I^n:=A
    \qquad
    (n\leq0).
\]
The preceding identifications give a canonical isomorphism
\[
\begin{aligned}
    \mathcal R_I^{\mathrm{ext}}(A)
    \otimes_RR'
    &\simeq
    \bigoplus_{n\in\mathbb Z}
    (I^n\otimes_RR')t^{-n}
    \\
    &\simeq
    \bigoplus_{n\in\mathbb Z}
    I'^n t^{-n}
    \\
    &=
    \mathcal R_{I'}^{\mathrm{ext}}(A').
\end{aligned}
\]
Taking relative spectra gives
\[
    D_{Z_0}X_0
    \times_{S_0}
    T_0
    \simeq
    D_{Z_T}X_T
\]
over \(\mathbb A^1_{T_0}\).

The closed subscheme \(Z_0\times\mathbb A^1\), the fibre over \(1\),
and the associated-graded fibre over \(0\) are all defined directly
from the extended Rees algebra.  The displayed isomorphism therefore
preserves them and the canonical identifications
\[
    (D_{Z_0}X_0)_1
    \simeq
    X_0,
    \qquad
    (D_{Z_0}X_0)_0
    \simeq
    N_{Z_0/X_0},
\]
and their base changes.  These affine identifications are compatible
with restriction and hence glue.
\end{proof}

\begin{theorem}[Functorial ordinary homotopy purity]
\label{thm:ordinary-homotopy-purity-functorial}
Let \(S_0\) be an ordinary quasi-compact quasi-separated scheme, and let
\[
    Z_0
    \hookrightarrow
    X_0
\]
be a closed immersion such that
\[
    X_0
    \longrightarrow
    S_0
\]
and
\[
    Z_0
    \longrightarrow
    S_0
\]
are smooth and of finite presentation.  There is a canonical
equivalence
\[
    \operatorname{pur}^{\mathrm{cl}}_{Z_0/X_0}
    \colon
    X_0/(X_0\setminus Z_0)
    \xrightarrow{\simeq}
    \operatorname{Th}^{\mathrm{cl}}_{S_0}
    \bigl(
        N_{Z_0/X_0}
    \bigr)
\]
in
\[
    \mathbf H_{\mathrm{cl},\bullet}(S_0).
\]

The equivalence is natural under arbitrary base change between
quasi-compact quasi-separated base schemes and under étale morphisms of
smooth closed pairs.
\end{theorem}

\begin{proof}
We first compare the ordinary motivic category used here with the
nonequivariant unstable motivic category occurring in the cited purity
theorem.

The ordinary smooth finite-presentation site used here and the
nonequivariant smooth site used in the cited theorem admit the same
Nisnevich-dense subsite of smooth affine schemes of finite presentation
over \(S_0\).  Every smooth scheme in either site admits a Nisnevich
covering by such smooth affine schemes, and fibre products of objects
of this affine subsite admit Nisnevich coverings by objects of the same
subsite.  The comparison lemma for a Nisnevich basis therefore
identifies the corresponding categories of Nisnevich sheaves.

It remains to compare the homotopy localizations.  Let
\[
    a
    \colon
    Y
    \longrightarrow
    X
\]
be a torsor under a vector bundle of rank \(n\).  Choose a Zariski
covering family
\[
    \{U_\alpha\to X\}_\alpha
\]
which trivializes the torsor.  On every iterated intersection
\[
    U_{\alpha_0\ldots\alpha_r},
\]
there is an equivalence
\[
    Y\times_XU_{\alpha_0\ldots\alpha_r}
    \simeq
    U_{\alpha_0\ldots\alpha_r}
    \times
    \mathbb A^n.
\]
If \(F\) is an ordinary Nisnevich sheaf which is
\(\mathbb A^1\)-invariant, repeated affine-line invariance gives
\[
    F(U_{\alpha_0\ldots\alpha_r})
    \xrightarrow{\simeq}
    F
    \bigl(
        Y\times_XU_{\alpha_0\ldots\alpha_r}
    \bigr).
\]
Taking totalizations of the two Čech diagrams and using Nisnevich
descent gives
\[
    F(X)
    \xrightarrow{\simeq}
    F(Y).
\]
Thus Nisnevich descent and \(\mathbb A^1\)-invariance imply invariance
under vector-bundle torsors.

Conversely, invariance under vector-bundle torsors, applied to the
trivial torsor under the additive group of the trivial line bundle,
\[
    X\times\mathbb A^1
    \longrightarrow
    X,
\]
implies \(\mathbb A^1\)-invariance.  Hence the two homotopy
localizations have the same local objects.  Together with the
Nisnevich-site comparison above, this identifies the nonequivariant
unstable motivic category used in the cited purity theorem with
\[
    \mathbf H_{\mathrm{cl}}(S_0).
\]

We may therefore apply the functorial ordinary homotopy-purity theorem
\cite[Theorem~3.23 and the discussion following it]
{HoyoisSixOperations}
in the category
\[
    \mathbf H_{\mathrm{cl},\bullet}(S_0).
\]

The closed immersion
\[
    Z_0
    \hookrightarrow
    X_0
\]
is regular because both \(Z_0\) and \(X_0\) are smooth over \(S_0\).
Let
\[
    D_{Z_0}X_0
\]
denote the deformation space of \(X_0\) along \(Z_0\).  It is smooth
and of finite presentation over
\[
    S_0\times\mathbb A^1.
\]
Moreover,
\[
    Z_0\times\mathbb A^1
    \hookrightarrow
    D_{Z_0}X_0
\]
is a closed immersion of finite presentation, and the pair
\[
    \bigl(
        D_{Z_0}X_0,
        Z_0\times\mathbb A^1
    \bigr)
\]
is a smooth closed pair over
\[
    S_0\times\mathbb A^1.
\]
Consequently, after composing with
\[
    S_0\times\mathbb A^1
    \longrightarrow
    S_0,
\]
both
\[
    D_{Z_0}X_0
\]
and its quasi-compact open complement
\[
    D_{Z_0}X_0
    \setminus
    \bigl(
        Z_0\times\mathbb A^1
    \bigr)
\]
are smooth schemes of finite presentation over \(S_0\).  Thus the
pointed quotient
\[
    D_{Z_0}X_0/
    \left(
        D_{Z_0}X_0
        \setminus
        \bigl(
            Z_0\times\mathbb A^1
        \bigr)
    \right)
\]
is an object of
\[
    \mathbf H_{\mathrm{cl},\bullet}(S_0).
\]

The fibre of the deformation space over \(1\) is canonically
equivalent to \(X_0\), while its fibre over \(0\) is canonically
equivalent to the normal bundle
\[
    N_{Z_0/X_0}.
\]
The two fibre inclusions induce a natural zigzag
\[
    X_0/(X_0\setminus Z_0)
    \xrightarrow{\alpha_1}
    D_{Z_0}X_0/
    \left(
        D_{Z_0}X_0
        \setminus
        \bigl(
            Z_0\times\mathbb A^1
        \bigr)
    \right)
    \xleftarrow{\alpha_0}
    \operatorname{Th}^{\mathrm{cl}}_{S_0}
    \bigl(
        N_{Z_0/X_0}
    \bigr).
\]
The ordinary homotopy-purity theorem asserts that both
\[
    \alpha_1
    \qquad\text{and}\qquad
    \alpha_0
\]
are equivalences in
\[
    \mathbf H_{\mathrm{cl},\bullet}(S_0).
\]
Define
\[
    \operatorname{pur}^{\mathrm{cl}}_{Z_0/X_0}
    :=
    \alpha_0^{-1}\alpha_1.
\]

Let
\[
    f
    \colon
    T_0
    \longrightarrow
    S_0
\]
be a morphism of quasi-compact quasi-separated schemes, and form the
base-changed pair
\[
    Z_T
    :=
    Z_0\times_{S_0}T_0
    \hookrightarrow
    X_T
    :=
    X_0\times_{S_0}T_0.
\]
Motivic inverse image preserves colimits and carries motives of smooth
schemes to motives of their geometric base changes.  It therefore
gives canonical equivalences
\[
    f^*
    \bigl(
        X_0/(X_0\setminus Z_0)
    \bigr)
    \simeq
    X_T/(X_T\setminus Z_T)
\]
and
\[
    f^*
    \operatorname{Th}^{\mathrm{cl}}_{S_0}
    \bigl(
        N_{Z_0/X_0}
    \bigr)
    \simeq
    \operatorname{Th}^{\mathrm{cl}}_{T_0}
    \bigl(
        N_{Z_T/X_T}
    \bigr).
\]

By
\Ref{lem:deformation-space-arbitrary-base-change},
there is a canonical equivalence
\[
    D_{Z_0}X_0
    \times_{S_0}
    T_0
    \simeq
    D_{Z_T}X_T
\]
over \(\mathbb A^1_{T_0}\), compatible with the closed immersions
\[
    Z_0\times\mathbb A^1
    \hookrightarrow
    D_{Z_0}X_0
\]
and
\[
    Z_T\times\mathbb A^1
    \hookrightarrow
    D_{Z_T}X_T,
\]
with the fibres over \(1\), and with the normal-bundle fibres over
\(0\).  Consequently, inverse image carries the deformation-space
zigzag defining
\[
    \operatorname{pur}^{\mathrm{cl}}_{Z_0/X_0}
\]
to the deformation-space zigzag defining
\[
    \operatorname{pur}^{\mathrm{cl}}_{Z_T/X_T}.
\]
Under the preceding source and target identifications,
\[
    f^*\alpha_1
\]
and
\[
    f^*\alpha_0
\]
identify with the two specialization morphisms for
\[
    Z_T
    \hookrightarrow
    X_T.
\]
It follows that
\[
    f^*
    \operatorname{pur}^{\mathrm{cl}}_{Z_0/X_0}
    \simeq
    \operatorname{pur}^{\mathrm{cl}}_{Z_T/X_T}.
\]

Now let
\[
    (X'_0,Z'_0)
    \longrightarrow
    (X_0,Z_0)
\]
be an étale morphism of smooth closed pairs over \(S_0\).  The assertion
is local on \(X_0\).  Write
\[
    X_0
    =
    \Spec(A),
    \qquad
    Z_0
    =
    \Spec(A/I),
\]
and, after restricting to an affine open of \(X'_0\), write
\[
    X'_0
    =
    \Spec(A').
\]
Since
\[
    A
    \longrightarrow
    A'
\]
is étale, it is flat.  The Cartesian condition on the closed pairs
identifies the ideal of \(Z'_0\) with
\[
    I'
    =
    IA'
    \simeq
    I\otimes_AA'.
\]
Flatness gives, for every \(n\geq0\), canonical isomorphisms
\[
    I^n\otimes_AA'
    \xrightarrow{\simeq}
    (I')^n.
\]
Consequently, there is a canonical isomorphism of extended Rees
algebras
\[
    \mathcal R_I^{\mathrm{ext}}(A)
    \otimes_A
    A'
    \xrightarrow{\simeq}
    \mathcal R_{I'}^{\mathrm{ext}}(A').
\]
Taking relative spectra gives a Cartesian morphism of deformation
spaces
\[
    D_{Z'_0}X'_0
    \longrightarrow
    D_{Z_0}X_0
\]
which is compatible with the parameter on \(\mathbb A^1\), the closed
subschemes
\[
    Z'_0\times\mathbb A^1
    \qquad\text{and}\qquad
    Z_0\times\mathbb A^1,
\]
and the fibres over \(1\) and \(0\).  On the zero fibre it is the
canonical pullback morphism of normal bundles.

These affine constructions are compatible with restriction and
therefore glue to a morphism between the two deformation-space
zigzags.  Naturality of
\[
    \alpha_1
    \qquad\text{and}\qquad
    \alpha_0
\]
then gives the asserted compatibility of the purity equivalences with
étale morphisms of smooth closed pairs.
\end{proof}

\begin{theorem}[Brave new homotopy purity]
\label{thm:brave-new-homotopy-purity}
Let
\[
    i
    \colon
    Z
    \hookrightarrow
    X
\]
be a smooth closed pair over \(S\).  There is a canonical equivalence
\[
    \operatorname{pur}_{Z/X}
    \colon
    X/(X\setminus Z)
    \xrightarrow{\simeq}
    \operatorname{Th}_S
    \bigl(
        \mathbb N_{Z/X}
    \bigr)
\]
in
\[
    \mathbf H_\bullet(S).
\]
This equivalence is compatible with arbitrary base change in \(S\) and
with étale morphisms of smooth closed pairs.
\end{theorem}

\begin{proof}
By
\Ref{prop:classical-comparison-quotient-thom},
the pointed comparison equivalence identifies the source and target
with
\[
    X_{\mathrm{cl}}/
    \bigl(
        X_{\mathrm{cl}}\setminus Z_{\mathrm{cl}}
    \bigr)
\]
and
\[
    \operatorname{Th}^{\mathrm{cl}}_{S_{\mathrm{cl}}}
    \bigl(
        N_{Z_{\mathrm{cl}}/X_{\mathrm{cl}}}
    \bigr),
\]
respectively.

The functorial ordinary purity theorem
\Ref{thm:ordinary-homotopy-purity-functorial}
supplies a canonical equivalence
\[
    \operatorname{pur}^{\mathrm{cl}}_{Z_{\mathrm{cl}}/X_{\mathrm{cl}}}
    \colon
    X_{\mathrm{cl}}/
    \bigl(
        X_{\mathrm{cl}}\setminus Z_{\mathrm{cl}}
    \bigr)
    \xrightarrow{\simeq}
    \operatorname{Th}^{\mathrm{cl}}_{S_{\mathrm{cl}}}
    \bigl(
        N_{Z_{\mathrm{cl}}/X_{\mathrm{cl}}}
    \bigr).
\]

Since
\[
    \Theta_{S,\bullet}
    \colon
    \mathbf H_\bullet(S)
    \xrightarrow{\simeq}
    \mathbf H_{\mathrm{cl},\bullet}(S_{\mathrm{cl}})
\]
is fully faithful, there is a unique morphism
\[
    \operatorname{pur}_{Z/X}
    \colon
    X/(X\setminus Z)
    \longrightarrow
    \operatorname{Th}_S
    \bigl(
        \mathbb N_{Z/X}
    \bigr)
\]
whose image under \(\Theta_{S,\bullet}\) is the ordinary purity
equivalence.  Since \(\Theta_{S,\bullet}\) reflects equivalences,
\(\operatorname{pur}_{Z/X}\) is an equivalence.

Let
\[
    f
    \colon
    T
    \longrightarrow
    S
\]
be arbitrary.  Base change of the closed pair gives
\[
    Z_T
    \hookrightarrow
    X_T.
\]
Cotangent base change gives
\[
    N^*_{Z_T/X_T}
    \simeq
    f_Z^*N^*_{Z/X},
\]
and therefore
\[
    \mathbb N_{Z_T/X_T}
    \simeq
    \mathbb N_{Z/X}\times_ST.
\]
Compatibility of \(\Theta\) with inverse image identifies the image of
\[
    f^*\operatorname{pur}_{Z/X}
\]
with the pullback of the ordinary purity equivalence.  By
\Ref{thm:ordinary-homotopy-purity-functorial},
this is the ordinary purity equivalence for
\[
    Z_T
    \hookrightarrow
    X_T.
\]
Full faithfulness of \(\Theta_{T,\bullet}\) therefore gives
\[
    f^*\operatorname{pur}_{Z/X}
    =
    \operatorname{pur}_{Z_T/X_T}.
\]

The same argument, using the étale naturality in
\Ref{thm:ordinary-homotopy-purity-functorial},
proves compatibility with étale morphisms of smooth closed pairs.
\end{proof}

\begin{corollary}[Purity for a zero section]
\label{cor:purity-zero-section}
Let \(E\) be finite locally free on a smooth spectral \(S\)-scheme
\(Z\), and let
\[
    s
    \colon
    Z
    \hookrightarrow
    \mathbb V_Z(E)
\]
be the zero section.  Then the purity equivalence identifies
\[
    \mathbb V_Z(E)/
    \bigl(
        \mathbb V_Z(E)\setminus Z
    \bigr)
\]
with
\[
    \operatorname{Th}_S(E).
\]
\end{corollary}

\begin{proof}
The conormal module of the zero section is canonically \(E\), by the
cotangent calculation for a relative free commutative algebra.
\end{proof}

\begin{proposition}[Multiplicativity of Thom spaces]
\label{prop:unstable-thom-multiplicativity}
Let \(E\) and \(F\) be finite locally free quasi-coherent modules on a
smooth spectral \(S\)-scheme \(p:Z\to S\).  There is a canonical
equivalence internal to \(Z\):
\[
    \operatorname{Th}_Z(E\oplus F)
    \simeq
    \operatorname{Th}_Z(E)
    \wedge_Z
    \operatorname{Th}_Z(F).
\]
It is compatible with base change.  Moreover,
\[
    \operatorname{Th}_S(E)
    \simeq
    p_{\sharp,\bullet}
    \operatorname{Th}_Z(E).
\]
\end{proposition}

\begin{proof}
Put
\[
    V
    :=
    \mathbb V_Z(E),
    \qquad
    W
    :=
    \mathbb V_Z(F),
\]
and write
\[
    V^\times
    :=
    V\setminus Z,
    \qquad
    W^\times
    :=
    W\setminus Z.
\]
There is a canonical equivalence of relative linear spaces over \(Z\):
\[
    \mathbb V_Z(E\oplus F)
    \simeq
    V\times_ZW.
\]
The complement \(C\) of the joint zero section in \(V\times_ZW\) is
the union
\[
    C
    =
    (V^\times\times_ZW)
    \cup
    (V\times_ZW^\times),
\]
whose intersection is
\[
    V^\times\times_ZW^\times.
\]
Consequently, there is a Zariski distinguished square
\[
\begin{tikzcd}[column sep=large, row sep=large]
    V^\times\times_ZW^\times
        \arrow[r]
        \arrow[d]
    &
    V^\times\times_ZW
        \arrow[d]
    \\
    V\times_ZW^\times
        \arrow[r]
    &
    C.
\end{tikzcd}
\]
Applying the motivic Yoneda functor gives a cocartesian square in
\(\mathbf H(Z)\).  Since motives preserve products, it identifies
\(\mathrm M_Z(C)\) with the pushout
\[
    \bigl(
        \mathrm M_Z(V^\times)
        \times
        \mathrm M_Z(W)
    \bigr)
    \amalg_{
        \mathrm M_Z(V^\times)
        \times
        \mathrm M_Z(W^\times)
    }
    \bigl(
        \mathrm M_Z(V)
        \times
        \mathrm M_Z(W^\times)
    \bigr).
\]

For morphisms
\[
    A_0
    \longrightarrow
    A,
    \qquad
    B_0
    \longrightarrow
    B
\]
in a presentable Cartesian category whose product preserves colimits
separately, distributivity of products over pushouts and the definition
of the smash product give a canonical equivalence
\[
    (A/A_0)
    \wedge
    (B/B_0)
    \simeq
    \frac{A\times B}{
        (A_0\times B)
        \amalg_{A_0\times B_0}
        (A\times B_0)
    }.
\]
Apply this with
\[
    A
    =
    \mathrm M_Z(V),
    \qquad
    A_0
    =
    \mathrm M_Z(V^\times),
\]
and
\[
    B
    =
    \mathrm M_Z(W),
    \qquad
    B_0
    =
    \mathrm M_Z(W^\times).
\]
The preceding Zariski-excision square identifies the denominator with
\(\mathrm M_Z(C)\).  Hence
\[
\begin{aligned}
    \operatorname{Th}_Z(E)
    \wedge_Z
    \operatorname{Th}_Z(F)
    &\simeq
    \frac{
        \mathrm M_Z(V\times_ZW)
    }{
        \mathrm M_Z(C)
    }
    \\
    &\simeq
    \operatorname{Th}_Z(E\oplus F).
\end{aligned}
\]
Every construction in this argument commutes with base change: relative
linear spaces, zero sections, open complements, the Zariski-excision
square, products, and pushouts.  Thus the multiplicativity equivalence
is compatible with base change.

The pointed smooth direct image \(p_{\sharp,\bullet}\) preserves
colimits and sends the motives of \(\mathbb V_Z(E)\) and its open
complement to the same schemes viewed over \(S\).  Applying it to the
pushout defining \(\operatorname{Th}_Z(E)\) gives
\[
    p_{\sharp,\bullet}\operatorname{Th}_Z(E)
    \simeq
    \operatorname{Th}_S(E).
\]
\end{proof}

\section{The localization and purity package}
\label{sec:localization-purity-package}


\begin{theorem}[Localization, purity, and nil-invariance]
\label{thm:localization-purity-nil-invariance-package}
Let \(S\) be a quasi-compact quasi-separated spectral scheme, and let
\[
    Z
    \xrightarrow{i}
    S
    \xleftarrow{j}
    U
\]
be a complementary closed--open pair with \(j\) quasi-compact.  Then
the following structures and properties are canonically available.
\begin{enumerate}[label=(\roman*)]
    \item For every \(F\in\mathbf H(S)\), the localization square
    \[
    \begin{tikzcd}[column sep=large, row sep=large]
        j_\sharp j^*F
            \arrow[r]
            \arrow[d]
        &
        F
            \arrow[d]
        \\
        j_\sharp(*_U)
            \arrow[r]
        &
        i_*i^*F
    \end{tikzcd}
    \]
    is cocartesian.

    \item For every \(F\in\mathbf H_\bullet(S)\), there are natural
    localization sequences
    \[
        j_\sharp j^*F
        \longrightarrow
        F
        \longrightarrow
        i_*i^*F
    \]
    and
    \[
        i_*i^!F
        \longrightarrow
        F
        \longrightarrow
        j_*j^*F.
    \]

    \item The functors \(i_*,j_\sharp,j_*\) are fully faithful, and
    \[
        i_*
        \colon
        \mathbf H_\bullet(Z)
        \xrightarrow{\simeq}
        \mathbf H_{\bullet,Z}(S)
    \]
    is an equivalence.

    \item Closed base change holds:
    for every Cartesian square
    \[
    \begin{tikzcd}
        Y \arrow[r, "k"] \arrow[d, "g"']
        &
        X \arrow[d, "f"]
        \\
        Z \arrow[r, "i"']
        &
        S
    \end{tikzcd}
    \]
    there is a canonical equivalence
    \[
        f^*i_*
        \simeq
        k_*g^*
    \]
    on unpointed and pointed motivic spaces.

    \item The pointed closed projection formula holds:
    for
    \[
        F\in\mathbf H_\bullet(Z),
        \qquad
        G\in\mathbf H_\bullet(S),
    \]
    there is a canonical equivalence
    \[
        i_*F\wedge_SG
        \simeq
        i_*
        \bigl(
            F\wedge_Zi^*G
        \bigr).
    \]

    \item Pointed smooth--closed exchange holds:
    for every Cartesian square
    \[
    \begin{tikzcd}
        Y \arrow[r, "k"] \arrow[d, "q"']
        &
        X \arrow[d, "p"]
        \\
        Z \arrow[r, "i"']
        &
        S
    \end{tikzcd}
    \]
    with \(p\) smooth and of finite presentation, there is a canonical
    equivalence
    \[
        p_{\sharp,\bullet}k_*
        \simeq
        i_*q_{\sharp,\bullet}
    \]
    of functors
    \[
        \mathbf H_\bullet(Y)
        \longrightarrow
        \mathbf H_\bullet(S).
    \]

    \item For every smooth closed pair
    \[
        Z
        \hookrightarrow
        X
    \]
    over \(S\), there is a canonical homotopy-purity equivalence
    \[
        X/(X\setminus Z)
        \simeq
        \operatorname{Th}_S(\mathbb N_{Z/X}).
    \]

    \item The theory is derived nilpotent invariant:
    \[
        \mathbf H(S)
        \simeq
        \mathbf H(S_{\mathrm{cl}}),
        \qquad
        \mathbf H_\bullet(S)
        \simeq
        \mathbf H_\bullet(S_{\mathrm{cl}}).
    \]

    \item There are natural classical comparison equivalences
    \[
        \mathbf H(S)
        \simeq
        \mathbf H_{\mathrm{cl}}(S_{\mathrm{cl}})
    \]
    and
    \[
        \mathbf H_\bullet(S)
        \simeq
        \mathbf H_{\mathrm{cl},\bullet}(S_{\mathrm{cl}}).
    \]
\end{enumerate}
\end{theorem}

\begin{proof}
Localization is
\Ref{thm:spectral-morel-voevodsky-localization}.
The pointed localization sequences are
\Ref{thm:pointed-localization-cofiber-sequence}
and
\Ref{thm:pointed-localization-fiber-sequence}.
Recollement and support are
\Ref{thm:closed-open-recollement}
and
\Ref{thm:closed-support-kashiwara}.
Closed base change, the pointed closed projection formula, and pointed
smooth--closed exchange are
\Ref{thm:unstable-closed-base-change},
\Ref{thm:unstable-closed-projection-formula}, and
\Ref{thm:unstable-smooth-closed-base-change}.
Homotopy purity is
\Ref{thm:brave-new-homotopy-purity}.
Derived nilpotent invariance is
\Ref{thm:derived-nilpotent-invariance},
and classical comparison is
\Ref{thm:classical-comparison-brave-new-motivic-spaces}.
\end{proof}

The constructions of the chapter may be summarized by the recollement
diagram
\[
    \mathbf H_\bullet(Z)
    \;
    \substack{
        \xrightarrow{\ i_*\ }
        \\
        \xleftarrow[\ i^*\ ]{}
        \\
        \xleftarrow[\ i^!\ ]{}
    }
    \;
    \mathbf H_\bullet(S)
    \;
    \substack{
        \xrightarrow{\ j^*\ }
        \\
        \xleftarrow[\ j_\sharp\ ]{}
        \\
        \xleftarrow[\ j_*\ ]{}
    }
    \;
    \mathbf H_\bullet(U),
\]
the purity equivalence
\[
    X/(X\setminus Z)
    \simeq
    \operatorname{Th}_S(\mathbb N_{Z/X}),
\]
and the truncation equivalence
\[
    \mathbf H(S)
    \simeq
    \mathbf H_{\mathrm{cl}}(S_{\mathrm{cl}}).
\]

The next chapter stabilizes
\[
    \mathbf H_\bullet(S)
\]
by Thom spaces.  The localization sequences then become exact triangles,
homotopy purity supplies the geometric Thom identifications, and the
closed exchange formulas become the first exceptional component of the
stable six-operation formalism.
